\chapter{Calabi--Yau Categories}
\label{ch:cy-categories}

A Calabi--Yau category is a dg category whose Serre functor is a pure degree shift. This single condition, due independently to Kontsevich and to Costello, organises the diverse sources of Calabi--Yau geometry (coherent sheaves, Fukaya categories, matrix factorisations, noncommutative resolutions) into a uniform categorical framework. It is the framework on which the Vol~III correspondence programme $\{\Phi_d\}_{d \geq 1}$ acts: each constructed $\Phi_d$ factors as $\Phi_d = \SpCh_{\Sigma_{d-1},C} \circ \PhiFA_d$ (Theorem~\ref{thm:phi-two-stage-factorisation}; Definition~\ref{def:phi-fa-and-sp}), with stage~$1$ $\PhiFA_d$ taking a cyclic $\Ainf$ Calabi--Yau category of dimension $d$ to an $E_d$-holomorphic factorisation algebra on the CY$_d$ variety, and stage~$2$ $\SpCh_{\Sigma_{d-1},C}$ specialising along a $(d-1)$-cycle and reference curve to produce an $\Etwo$-chiral algebra at $d = 2$ or an $\Eone$-chiral algebra on the witnessed $d \geq 3$ loci; the native operadic level $n_{\mathrm{native}}(d)$ on $C$ (Proposition~\ref{prop:native-en-level}) is part of the dimension-specific data. This chapter fixes the categorical input that enters stage~$1$, recording the definitions, the Hochschild structure, the cyclic $\Ainf$ enhancement, and the interface to $\{\Phi_d\}$ (Chapter~\ref{ch:cy-to-chiral}).

\begin{remark}[Chapter scope and cross-volume anchors]
\label{rem:cy-cat-cross-volume-anchors}
The categorical input fixed by this chapter feeds the correspondence programme $\{\Phi_d\}$ of Theorem~\textup{\ref{thm:phi-platonic}}, characterised by the four universal properties \textup{(U1)}--\textup{(U4)}. The downstream evaluations of $\{\Phi_d\}$ supplied by the cyclic $\Ainf$ data of this chapter are the dimension-stratified spectrum of $\kappa_{\mathrm{ch}}$ values catalogued in Theorem~\textup{\ref{thm:kappa-stratification-by-d}} (Chapter~\textup{\ref{ch:cy-d-kappa-stratification}}) and the Borcherds-weight invariant $\kappa_{\mathrm{BKM}} = c_N(0)/2$ of Proposition~\textup{\ref{prop:bkm-weight-universal}} (Chapter~\textup{\ref{ch:k3-times-e}}). On the Vol~I side, the categorical Hochschild data of Section~\textup{\ref{sec:cy-hochschild}} is the cohomological shadow of the chiral universal Maurer--Cartan element of Vol~I Chapter~\textup{\texttt{chapters/theory/chiral\_climax\_platonic.tex} (Vol~I)}, and the Koszul conductor identity $K = -c_{\mathrm{ghost}}(\BRST)$ of Vol~I Chapter~\textup{\texttt{chapters/theory/universal\_conductor\_K\_platonic.tex} (Vol~I)} fixes the family-dependent constants entering the kappa-spectrum.
\end{remark}

\begin{remark}[Chain-level and \texorpdfstring{$(\infty,1)$}{(inf,1)}-categorical: two load-bearing lenses]
\label{rem:cy-cat-chain-infinity-parity}
Both lenses are load-bearing throughout the chapter and neither
replaces the other. The chain-level lane supplies the explicit
cyclic $\Ainf$ enhancement (Theorem~\ref{thm:cyclic-ainf-enhancement}),
the explicit Hochschild complexes of Section~\ref{sec:cy-hochschild},
and the ambient-qualified Serre-duality computations driving the
$K3$, $K3\times E$, and local $\mathbb P^2$ examples. The
$(\infty,1)$-categorical lane supplies the definition of
$\CY_d\text{-}\Cat$ as an $\infty$-category, the Maulik--Okounkov /
Kontsevich--Soibelman stable-envelope functoriality, and the
conditional $E_3$ obstruction package used for
$\Phi_3^{(\Sigma_2,C)}$ on the H1--H4 locus
(Theorem~\ref{thm:cy-to-chiral-d3}). In that package
$\HH^{-2}_{E_1}$ vanishing is a theorem only after the comparison map
to the obstruction complex and the complete/exhaustive/separated
strongly convergent filtration with empty total-degree $-2$ line have
been supplied, or after the corrected TCFT carrier gives the same
comparison. Every theorem is stated in the
lane in which its proof works, with scope qualifiers
(ambient, framing, $\Ainf$-level) separating the two readings;
neither lane is ``the real theorem'' of which the other is a shadow.
\end{remark}

\begin{remark}[Convention: \texorpdfstring{$\kappa_\bullet$}{kappa-.} subscripts in Vol~III]
\label{rem:cy-cat-kappa-convention}
Throughout this chapter and the rest of Vol~III the symbol $\kappa_\bullet$ ranges over the closed indexing set $\bullet \in \{\mathrm{ch}, \mathrm{cat}, \mathrm{BKM}, \mathrm{fiber}\}$: every occurrence carries one of these four subscripts, and each subscript names a distinct invariant. Indexed against the two-stage factorisation $\Phi_d = \SpCh_{\Sigma_{d-1},C} \circ \PhiFA_d$ (Theorem~\ref{thm:phi-two-stage-factorisation}), the four entries split into three tiers:
\begin{itemize}
 \item \emph{Tier~(i), CY-datum intrinsic} (independent of stage~$1$ or stage~$2$): $\kappa_{\mathrm{cat}} = \chi(\cO_X)$, the holomorphic Euler characteristic of the CY$_d$ variety. Read off the Hodge diamond alone.
 \item \emph{Tier~(ii), stage~$1$ invariant} (determined by $\PhiFA_d(\cC)$ on $X$ before any specialisation): $\kappa_{\mathrm{fiber}} = \mathrm{rank}(\Lambda)$, the rank of the lattice of local / divisor classes on $X$ seen by the $E_d$-factorisation algebra.
 \item \emph{Tier~(iii), stage~$2$ specialisation-dependent} (determined by the stage-$2$ output $\SpCh_{\Sigma_{d-1},C}(\PhiFA_d(\cC))$ on the reference curve): $\kappa_{\mathrm{ch}}(\Phi_d(\cC))$ (chiral modular characteristic as the Hodge-supertrace reading on $C$) and the Borcherds-weight $\kappa_{\mathrm{BKM}}(\Phi_{N}(\cC))$ (weight of the BKM denominator automorphic form attached to a specialisation cycle).
\end{itemize}
The statement ``algebraizations share $\kappa_{\mathrm{cat}}$'' is vacuous: $\kappa_{\mathrm{cat}}$ is a tier~(i) invariant of $X$, not a property of $\cC$ or of a specialisation.
\end{remark}

\section{Smooth and proper dg categories}
\label{sec:smooth-proper-dg}

The Serre functor required by the CY condition exists only for categories
that are both compact over themselves and finitely presented in their
$\Hom$-complexes. These are the smooth and proper conditions.

\begin{definition}[Smooth dg category]
\label{def:smooth-dg-cat}
A dg category $\cC$ over $k$ is \emph{smooth} if the diagonal bimodule $\cC \in \cC^{\op} \otimes \cC\text{-}\Mod$ is a perfect bimodule (compact in the derived category of bimodules).
\end{definition}

\begin{definition}[Proper dg category]
\label{def:proper-dg-cat}
A dg category $\cC$ is \emph{proper} if for all objects $X, Y \in \cC$, the Hom-complex $\Hom_\cC(X,Y)$ has finite-dimensional total cohomology.
\end{definition}

A category that is both smooth and proper is called \emph{saturated}. Saturated categories admit a Serre functor $S_\cC$, a distinguished autoequivalence characterised by the natural isomorphism $\Hom_\cC(X, Y) \simeq \Hom_\cC(Y, S_\cC X)^\vee$. Classical Serre duality recovers $S_{D^b(\Coh(X))} = (- \otimes \omega_X)[\dim X]$; the Calabi--Yau condition asks that the twist $\omega_X$ be trivial, leaving only the shift.

\section{The Calabi--Yau condition}
\label{sec:cy-condition}

A saturated category carries a Serre functor $S_\cC = (-)\otimes\omega_\cC[d]$.
Two pieces can be trivial: the twist $\omega_\cC$ and the shift $[d]$. Asking
$\omega_\cC$ to be trivial leaves a pure shift and is the Calabi--Yau condition.

\begin{definition}[Calabi--Yau category, after Kontsevich]
\label{def:cy-category}
A smooth dg (or $\Ainf$) category $\cC$ is \emph{Calabi--Yau of dimension $d$} if the Serre functor is a pure shift: $S_\cC \simeq [d]$. Equivalently, there is a quasi-isomorphism of $\cC$-bimodules
\[
 \cC \xrightarrow{\;\sim\;} \cC^\vee[d]
\]
where $\cC^\vee = \RHom_{\cC^{\op} \otimes \cC}(\cC, \cC^{\op} \otimes \cC)$ is the inverse dualizing bimodule. Equivalently again, there exists a non-degenerate cyclic pairing
\[
 \langle \cdot, \cdot \rangle_d \colon \cC(a, b) \otimes \cC(b, a) \longrightarrow k[-d]
\]
of degree $-d$, satisfying the graded symmetry $\langle f, g \rangle_d = (-1)^{|f|\,|g|} \langle g, f \rangle_d$.
\end{definition}

The equivalence of the bimodule formulation with the cyclic pairing formulation is due to Costello~\cite{Costello2005}; see also Kontsevich's Mirror Symmetry address~\cite{Kontsevich1994} for the original categorical statement and Keller~\cite{Keller2006} for the dg-categorical reformulation used here.

\begin{remark}
\label{rem:cy-vs-smooth-cy}
The CY condition combines smoothness (the bimodule $\cC$ is compact) with the Serre functor being a shift. A smooth proper CY category is \emph{saturated} in the sense of Kontsevich--Soibelman. Non-proper ``relative'' CY categories also appear (matrix factorisations, Ginzburg dg algebras); they require the Hochschild invariants of Section~\ref{sec:cy-hochschild} to be interpreted in a compactly-supported or negative-cyclic refinement.
\end{remark}

\begin{example}[Fukaya category]
\label{ex:fukaya-cy}
For a compact symplectic manifold $(M, \omega)$ with $c_1(M) = 0$, the Fukaya category $\Fuk(M)$ is CY of dimension $\dim_{\mathbb{R}}(M)/2 = n$. The CY pairing arises from the cyclic structure on the $\Ainf$-algebra of Lagrangian Floer cochains (Fukaya; Seidel~\cite{Seidel2008}).
\end{example}

\begin{example}[Derived category of a CY manifold]
\label{ex:derived-cy}
For a smooth projective CY manifold $X$ of complex dimension $n$ (i.e.\ $\omega_X \simeq \cO_X$ and $H^i(X, \cO_X) = 0$ for $0 < i < n$), the bounded derived category $D^b(\Coh(X))$ is CY of dimension $n$. The CY pairing is Serre duality composed with the $\cO_X$-trivialisation.
\end{example}

\begin{example}[Matrix factorizations]
\label{ex:mf-cy}
For $W \colon \mathbb{A}^n \to \mathbb{A}^1$ a polynomial with isolated singularity, the category $\MF(W)$ of matrix factorizations is CY of dimension $n - 2$ (Dyckerhoff; Orlov). For the ADE singularities in two variables ($n = 2$), $\MF(W)$ is $\CY_0$ (semisimple, with the structure of a Frobenius algebra). For threefold singularities ($n = 4$, e.g.\ $W = x_1^2 + \cdots + x_4^2$), $\MF(W)$ is $\CY_2$. The exponent $n - 2$ (not $n - 1$) records the projectivisation of the $\mathbb{G}_m$-action on $W^{-1}(0)$.
\end{example}

\section{Hochschild cohomology and the CY structure}
\label{sec:cy-hochschild}

Hochschild cohomology is the universal invariant of a dg category. For a CY$_d$ category, Serre duality forces additional structure: an $E_2$-algebra structure on $\HH^\bullet(\cC)$ (Kontsevich--Soibelman, Tamarkin) and, in the presence of the CY pairing, a $(d-2)$-shifted Poisson / BV structure.

\begin{remark}[Three Hochschild theories (chiral / topological / categorical): NEVER conflate]
\label{rem:cy-cat-three-hochschild}
The geometry determines which Hochschild theory is at play. Three theories coexist in this programme; they are distinct functors on distinct sources, with distinct outputs:
\begin{enumerate}
 \item \emph{Categorical Hochschild} $\HH^\bullet_{\mathrm{cat}}$ (this chapter): input a dg category $\cC$; output $E_2$-algebra (Deligne) carrying a $(2{-}d)$-shifted Poisson / $\BV$ structure when $\cC$ is $\CY_d$. Source: dg categories. The Vol~III categorical kappa is $\kappa_{\mathrm{cat}}(\cC) = \chi(\cO_X)$ when $\cC = D^b(\Coh(X))$.
 \item \emph{Topological Hochschild} $\HH^\bullet_{\mathrm{top}}$: input an $E_1$-algebra; output $E_2$-algebra (Deligne; Lurie). Source: associative algebras over a base ring. This is the Hochschild theory operative in Vol~II for the slab algebra of 3d HT theories.
 \item \emph{Chiral Hochschild} $\ChirHoch^\bullet$: input an $E_\infty$-chiral algebra on a curve $X$; output concentrated in cohomological degrees $\{0,1,2\}$ on the curve-ambient reading (Vol~I Theorem~H; $d=1$ Ran ambient). Under $\Phi_d$ with $d \ge 2$, the concentration enlarges to $\{0,1,2,d\}$ in the cohomological lane; chain-level is infinite for class-M inputs at $d \ge 3$. Source: chiral algebras. Carries a chiral Gerstenhaber bracket (from OPE residues on $\mathrm{FM}_k(C)$) distinct from the topological Gerstenhaber bracket (from the little $2$-disks operad); the two agree only for $E_\infty$-inputs via formality.
\end{enumerate}
The CY-to-chiral correspondence programme $\{\Phi_d\}$ of Chapter~\textup{\ref{ch:cy-to-chiral}} maps the categorical Gerstenhaber bracket on $\HH^\bullet_{\mathrm{cat}}(\cC)$ to the chiral convolution bracket on $\ChirHoch^\bullet(\Phi_d(\cC))$; the two are distinct objects on distinct sources. We therefore qualify every Hochschild statement that crosses this bridge with one of $\{\mathrm{cat},\mathrm{top},\mathrm{chiral}\}$, naming the theory under consideration.
\end{remark}

\begin{definition}[Hochschild cohomology]
\label{def:hochschild-cohom}
The Hochschild cohomology of $\cC$ is
\[
 \HH^\bullet(\cC) \;=\; \RHom_{\cC^{\op} \otimes \cC}(\cC, \cC),
\]
the derived endomorphisms of the diagonal bimodule. It is a graded-commutative algebra under Yoneda product and carries the Gerstenhaber bracket of degree $-1$. Throughout this chapter, $\HH^\bullet$ without qualifier means $\HH^\bullet_{\mathrm{cat}}$ in the sense of Remark~\textup{\ref{rem:cy-cat-three-hochschild}}.
\end{definition}

\begin{theorem}[Deligne conjecture; Kontsevich--Soibelman, Tamarkin]
\label{thm:deligne-conjecture}
The pair (Yoneda product, Gerstenhaber bracket) on $\HH^\bullet(\cC)$ is the homology of an $E_2$-algebra structure fixed after choosing the Kontsevich--Tamarkin formality torsor point.
\end{theorem}

For a CY$_d$ category, the cyclic pairing relates $\HH^\bullet(\cC)$ to $\HH_\bullet(\cC)[d]$. Transporting structure through this duality produces the higher Calabi--Yau brackets.

\begin{proposition}[CY bracket spectrum]
\label{prop:cy-bracket-spectrum}
For a $\CY_d$ category $\cC$, the Hochschild invariants carry the following additional structure:
\begin{itemize}
 \item $d = 1$: a BV operator $\Delta \colon \HH^\bullet(\cC) \to \HH^{\bullet - 1}(\cC)$ whose deviation from being a derivation recovers the Gerstenhaber bracket;
 \item $d = 2$: a degree $-2$ Poisson bracket $\{-, -\}_2$ on $\HH^\bullet(\cC)$, part of the $2$-shifted symplectic structure on the moduli of objects (PTVV shift law: Pantev--To\"en--Vaqui\'e--Vezzosi~\cite{PTVV2013}, $d$-shifted symplectic on $\mathrm{Map}(X, \mathrm{Pt}/G)$ for CY$_d$ target);
 \item $d \geq 3$: a $(d - 2)$-shifted Poisson structure / $P_{d-1}$-algebra (PTVV, loc.\ cit.), refining to the Lie-cobar level in the cyclic setting.
\end{itemize}
\end{proposition}

\begin{theorem}[HKR for smooth varieties]
\label{thm:hkr-smooth}
For $X$ a smooth projective variety over $k$ of characteristic zero, the Hochschild--Kostant--Rosenberg isomorphism gives
\[
 \HH^p(D^b(\Coh(X))) \;\simeq\; \bigoplus_{q + r = p} H^q\big(X, \wedge^r T_X\big).
\]
When $X$ is CY$_n$, the trivialisation $\omega_X \simeq \cO_X$ identifies $\wedge^r T_X \simeq \Omega_X^{n - r}$, so
\[
 \HH^p(D^b(\Coh(X))) \;\simeq\; \bigoplus_{q + r = p} H^q\big(X, \Omega_X^{n - r}\big).
\]
The sign and compatibility with the Mukai pairing are recorded in C{\u a}ld{\u a}raru~\cite{Caldararu2005}.
\end{theorem}

\begin{example}[Hochschild cohomology of \texorpdfstring{$K3$}{K3}]
\label{ex:hh-k3}
For a $K3$ surface $X$, the Hodge diamond is
\[
 h^{0,0} = 1,\; h^{2,0} = h^{0,2} = 1,\; h^{1,1} = 20,\; h^{2,2} = 1,
\]
with all other entries zero, and $H^q(\wedge^r T_X) \simeq H^q(\Omega_X^{2-r})$ by the CY$_2$ trivialisation $\omega_X \simeq \cO_X$. Applied to the Kontsevich HKR sum
$\HH^p(D^b(\Coh(X))) = \bigoplus_{q + r = p} H^q(X, \wedge^r T_X)$ (Theorem~\ref{thm:hkr-smooth}), the single non-vanishing summands at each total degree are
\[
 \HH^0 \simeq H^0(\cO_X) \simeq k, \qquad
 \HH^1 \simeq 0,
\]
\[
 \HH^2 \simeq H^0(\wedge^2 T_X) \,\oplus\, H^1(T_X) \,\oplus\, H^2(\cO_X) \simeq k \oplus k^{20} \oplus k \simeq k^{22},
\]
\[
 \HH^3 \simeq 0, \qquad
 \HH^4 \simeq H^2(\wedge^2 T_X) \simeq k,
\]
with the $H^1(T_X) \simeq k^{20}$ summand recording the infinitesimal K3 moduli and $H^0(\wedge^2 T_X) \simeq H^2(\cO_X) \simeq k$ arising from the trivialisation of $\wedge^2 T_X \simeq \cO_X$. The total dimension is
\[
 \dim_k \HH^\bullet(D^b(\Coh(K3))) \;=\; 1 + 0 + 22 + 0 + 1 \;=\; 24,
\]
matching the topological Euler characteristic $\chi(K3) = 24$. Collapsing by cohomological row rather than total degree (the Mukai grading of C{\u a}ld{\u a}raru) instead splits this $24$ as $(1{+}1) + 20 + (1{+}1) = 2 + 20 + 2$ across the three even cohomological rows of $K3$; this is the grading used downstream in the Mukai-lattice picture, distinct from the Kontsevich HKR grading above. The Mukai pairing on $\HH^\bullet$ is the non-degenerate pairing $\langle -, - \rangle_\mathrm{Mukai}$ of Mukai~\cite{Mukai1987}; it coincides with the CY pairing produced by Definition~\ref{def:cy-category} at $d = 2$.
\end{example}

\begin{remark}[\texorpdfstring{$\kappa_\mathrm{cat}$}{kappa-cat} as Hodge supertrace; comparison with \texorpdfstring{$\kappa_{\mathrm{ch}}$}{kappa-ch}]
\label{rem:kappa-cat-from-chi}
For $X$ a compact CY$_d$, the categorical kappa of the Vol~III spectrum is the Hodge-supertrace formula
\[
 \kappa_{\mathrm{cat}}(X) \;=\; \chi(\cO_X) \;=\; \sum_{q=0}^{d} (-1)^q\, h^{0,q}(X),
\]
the holomorphic Euler characteristic. For K3: $\chi(\cO_{K3}) = 1 - 0 + 1 = 2$, so $\kappa_{\mathrm{cat}}(K3) = 2$. At every \emph{odd} dimension $d$, Serre duality forces $h^{0,q}(X) = h^{0,d-q}(X)$, and the supertrace pair-cancels, giving
\[
 \chi(\cO_X) \;=\; 0 \qquad (\text{compact CY}_d,\ d\ \text{odd; Vol~III \texttt{prop:chi-O-vanishes-odd-d})}.
\]
The chiral kappa $\kappa_{\mathrm{ch}}(\Phi_d(\cC))$ is the compact Hodge-supertrace branch unless a construction superscript is displayed; the equality $\kappa_{\mathrm{ch}} = \kappa_{\mathrm{cat}}$ holds at $d = 2$ with $h^{1,0}(X) = 0$ (Theorem~\textup{\ref{thm:cy-d-d2}}; Proposition~\textup{\ref{prop:cy-kappa-d2}}). At $d = 3$ the Stage-2 Heisenberg branch can be nonzero while the compact trace vanishes: $\kappa_{\mathrm{ch}}^{\mathrm{Heis}}(K3 \times E) = 3$, but $\kappa_{\mathrm{ch}}(A_{K3\times E})=\kappa_{\mathrm{cat}}(K3\times E)=0$ (Remark~\ref{rem:beauville-kappa-formula-subscript-split}). The dimension-stratified spectrum of $\kappa_{\mathrm{ch}}$ values is catalogued in Theorem~\textup{\ref{thm:kappa-stratification-by-d}} and prescribes the correct invariant in each $d$.

Numerical separation for K3: $\dim \HH^\bullet(D^b(\Coh(K3))) = 24$ (the topological Euler characteristic); $\kappa_{\mathrm{cat}}(K3) = \chi(\cO_{K3}) = 2$; $\kappa_{\mathrm{ch}}(\Phi(D^b(\Coh(K3)))) = 2$ ($d = 2$, $h^{1,0}=0$). For the $K3 \times E$ product: $\kappa_{\mathrm{cat}}(K3 \times E) = \chi(\cO_{K3}) \cdot \chi(\cO_E) = 2 \cdot 0 = 0$ by K\"unneth and $\kappa_{\mathrm{ch}}(A_{K3\times E})=0$ by compact Hodge supertrace, while $\kappa_{\mathrm{ch}}^{\mathrm{Heis}}(K3 \times E) = 2 + 1 = 3$ is the canonical Stage-2 Heisenberg--Cartan rank, Beauville-additive under tensor products (Vol~I Theorem~D; Remark~\ref{rem:beauville-kappa-formula-subscript-split}), and $\kappa_{\mathrm{BKM}}(K3 \times E) = c_1(0)/2 = 5$ by the Borcherds weight theorem (Proposition~\textup{\ref{prop:bkm-weight-universal}}, Vol~III); the Mukai lattice rank is $\kappa_{\mathrm{fiber}}(K3) = 24$. The four construction values $\{\kappa_{\mathrm{cat}}, \kappa_{\mathrm{ch}}^{\mathrm{Heis}}, \kappa_{\mathrm{BKM}}, \kappa_{\mathrm{fiber}}\} = \{0, 3, 5, 24\}$ for the K3 $\times$ E example are distinct invariants, not one invariant evaluated four times.
\end{remark}

\section{Cyclic \texorpdfstring{$\Ainf$}{A-inf} structure}
\label{sec:cyclic-ainf-cy}

The CY pairing in Definition~\ref{def:cy-category} is cohomological. The categorical input to $\Phi$ requires a chain-level enhancement: a cyclic $\Ainf$-structure for which every higher operation $m_n$ is compatible with the pairing.

\begin{definition}[Cyclic \texorpdfstring{$\Ainf$}{A-inf} category]
\label{def:cyclic-ainf-cat}
A \emph{cyclic $\Ainf$-structure of degree $-d$} on an $\Ainf$-category $\cC$ is a collection of pairings
\[
 \langle -, - \rangle_d \colon \cC(a_0, a_1) \otimes \cC(a_1, a_0) \longrightarrow k[-d]
\]
together with the cyclic symmetry condition
\[
 \langle m_n(f_1, \ldots, f_n), f_0 \rangle_d \;=\; (-1)^{\epsilon} \langle m_n(f_0, f_1, \ldots, f_{n-1}), f_n \rangle_d
\]
for all $n \geq 2$, where $\epsilon$ is the Koszul sign determined by the degrees of the $f_i$.
\end{definition}

\begin{definition}[Negative cyclic CY class]
\label{def:negcyc-cy}
A \emph{$d$-dimensional CY structure} on a smooth dg category $\cC$ is a class
\[
 [\sigma] \in \HC^-_{d}(\cC)
\]
in the negative cyclic homology, whose image under the canonical map $\HC^- \to \HH$ is a non-degenerate trace $\Tr \colon \HH_d(\cC) \to k$. The lift to $\HC^-$ is essential for the $\bS^d$-framing used in Chapter~\ref{ch:cy-to-chiral}.
\end{definition}

\begin{theorem}[Cyclic \texorpdfstring{$\Ainf$}{A-inf} enhancement]
\label{thm:cyclic-ainf-enhancement}
Every smooth proper CY$_d$ category admits a cyclic $\Ainf$-enhancement, unique up to $\Ainf$-quasi-isomorphism preserving the pairing. The proof is due to Kontsevich--Soibelman in the formal case, Costello~\cite{Costello2005} for the TCFT enhancement, and Seidel~\cite{Seidel2008} for the Fukaya case.
\end{theorem}

At $d = 2$ the enhancement is explicit enough to compute directly; it drives the $K3$ computations feeding the Borcherds denominator (Vol~III Chapter~\ref{ch:k3-times-e}). At $d = 3$ Sheridan~\cite{Sheridan2015} constructs the required cyclic enhancement for the quintic HMS model. The general theorem is conditional on the chain-level $\bS^3$-framing datum and on a fixed admissible specialisation $(\Sigma_2,C)$ (Theorem~\ref{thm:cy-to-chiral-d3}); the framing is automatic only on verified toric or formal loci.

\section{Interface with the CY-to-chiral correspondence programme \texorpdfstring{$\{\Phi_d\}$}{Phi-d}}
\label{sec:cy-cat-phi-interface}

The proved Vol~III functor is
\[
 \Phi \colon \CY_2\text{-}\Cat^{\mathrm{cyc}} \longrightarrow \Etwo\text{-}\mathrm{ChirAlg}.
\]
The $d = 3$ extension exists as an $\Eone$-chiral object-level
assignment on the witnessed locus, and the $\Eone$ output level is
rigid in the sense of Theorem~CY-A$_3$
($E_2$-lifting contractible on the shadow via
Theorem~\ref{thm:derived-framing-obstruction}), \emph{not} an
equivalence of categories with anything:
\[
 \Phi_3^{(\Sigma_2,C)}(\cC)
 :=
 \SpCh_{\Sigma_2,C}(\PhiFA_3(\cC))
 \in \Eone\text{-}\mathrm{ChirAlg}(C).
\]
This is an assignment on objects after H1--H4 and the admissible specialisation datum have been fixed, not a functor on all smooth proper CY$_3$ categories. Both readings use as input the data of this chapter: a smooth proper cyclic $\Ainf$ category with its negative cyclic CY class $[\sigma] \in \HC^-_d(\cC)$. Each constructed $\Phi_d$ factors through a canonical two-stage construction (Theorem~\ref{thm:phi-two-stage-factorisation}): Stage~1 assembles the native $E_d$-holomorphic factorisation algebra $\PhiFA_d(\cC)$ on the CY$_d$ variety from the Gerstenhaber bracket of degree $1-d$ on $\HH^\bullet(\cC)$ via Kontsevich--Tamarkin formality and Costello--Gwilliam--Li locality; Stage~2 applies the chiral specialisation functor $\SpCh_{\Sigma_{d-1},C}$ to a $(d-1)$-dimensional cycle and reference curve. The operadic level $n_{\mathrm{native}}(d)$ of the chiral output on $C$ is determined by the Gerstenhaber bracket degree alone (Proposition~\ref{prop:native-en-level}): $\infty$ at $d=1$, $2$ at $d=2$, and $1$ at $d \geq 3$. The construction (Chapter~\ref{ch:cy-to-chiral}) proceeds concretely by (i) extracting the Gerstenhaber algebra $\HH^\bullet(\cC)$, (ii) promoting it to a Lie conformal algebra using the CY pairing, (iii) taking the factorization envelope on a curve, and (iv) at $d = 2$ enhancing to $\Etwo$ using the $\bS^2$-framing.

\begin{theorem}[Categorical kappa equals chiral kappa at \texorpdfstring{$d = 2$}{d = 2}, \texorpdfstring{$h^{1,0} = 0$}{h-1,0 = 0}: CY-D\texorpdfstring{$_2$}{-2}]
\label{thm:cy-d-d2}
For $\cC$ a smooth proper $\CY_2$ category such that $\Phi(\cC) = \SpCh_{\Sigma_1,C}(\PhiFA_2(\cC))$ exists \emph{and} $\HH_{-1}(\cC) = 0$ (equivalently, in the geometric case $\cC = D^b(\Coh(X))$, the hypothesis $h^{1,0}(X) = 0$), the chiral modular characteristic, namely the tier-(iii) invariant read off the stage-$2$ output on $C$, satisfies
\[
 \kappa_{\mathrm{ch}}\big(\Phi(\cC)\big) \;=\; \chi^{\CY}(\cC) \;=\; \chi(\cO_X) \;=\; \kappa_{\mathrm{cat}}(X).
\]
The middle equality is the categorical Euler characteristic via the Mukai pairing (tier-(i) invariant of the input); the right equality is the Hodge-supertrace formula of Remark~\textup{\ref{rem:kappa-cat-from-chi}}. The identity is a coincidence of a tier-(iii) output with a tier-(i) input, specific to $d = 2$ with $h^{1,0}=0$: for $\cC = D^b(\Coh(X))$ with $X$ a K3 surface, $\kappa_{\mathrm{ch}}(\Phi(D^b(\Coh(K3)))) = \kappa_{\mathrm{cat}}(K3) = 2$. The canonical Stage-2 product chiral rank $\kappa_{\mathrm{ch}}^{\mathrm{Heis}}(K3 \times E) = 3$ obtains by Beauville additivity under tensor products; the compact total-space trace is $\kappa_{\mathrm{ch}}(A_{K3\times E})=0$. The four-way construction separation $\{\kappa_{\mathrm{cat}}, \kappa_{\mathrm{ch}}^{\mathrm{Heis}}, \kappa_{\mathrm{BKM}}, \kappa_{\mathrm{fiber}}\} = \{0, 3, 5, 24\}$ for K3 $\times$ E is laid out in Remark~\textup{\ref{rem:kappa-cat-from-chi}}.
\end{theorem}

\begin{remark}[Hypothesis \texorpdfstring{$h^{1,0} = 0$}{h-1,0 = 0} is essential at \texorpdfstring{$d = 2$}{d = 2}]
\label{rem:cy-d-d2-hypothesis}
The identification $\kappa_{\mathrm{ch}} = \kappa_{\mathrm{cat}}$ fails without the hypothesis $h^{1,0}(X) = 0$. Counterexample: the abelian surface $T^4$ has $h^{1,0}(T^4) = 2$ and $\chi(\cO_{T^4}) = 0$, so $\kappa_{\mathrm{cat}}(T^4) = 0$, but additivity gives $\kappa_{\mathrm{ch}}(\Phi(T^4)) = \kappa_{\mathrm{ch}}(\Phi(E)) + \kappa_{\mathrm{ch}}(\Phi(E)) = 1 + 1 = 2 \neq 0$. The sharper Beauville formula of Proposition~\textup{\ref{prop:beauville-kappa-formula}}, $\kappa_{\mathrm{ch}}(\Phi(D^b(\Coh(X)))) = \chi(\cO_X) + h^{1,0}(X)$, reduces to $\kappa_{\mathrm{cat}} = \chi(\cO_X)$ exactly in the simply-connected $h^{1,0} = 0$ case. The most common silent failure mode is applying $\kappa_{\mathrm{ch}} = \chi(\cO_X)$ at odd $d$ where Serre forces $\chi(\cO_X) = 0$ but $\kappa_{\mathrm{ch}}$ is generically nonzero.
\end{remark}

\begin{remark}[The \texorpdfstring{$d = 3$}{d = 3} case]
\label{rem:cy-cat-d3}
For $d = 3$, the assignment $\Phi_3^{(\Sigma_2,C)}$ is constructed on the stated object-level locus only after the H1--H4 data, the admissible specialisation, and the obstruction comparison data of Theorem~\ref{thm:cy-to-chiral-d3} have been fixed. The required input is not the slogan $\HH^{-2}_{E_1}=0$: it is either a corrected TCFT carrier $B^{(2)}_{\mathrm{TCFT}}$ with its comparison to the obstruction complex, or an explicit $\HH^{-2}_{E_1}$ filtration theorem with complete/exhaustive/separated strong convergence and empty total-degree $-2$ line. The raw carrier $B^{(2)}_{\mathrm{term}}$ is excluded by the strict witness $[m_3,B^{(2)}_{\mathrm{term}}][a|a|a|a|b]=2\alpha[b]\neq0$. The $\mathbb{C}^3$ Hall/representation chain is verified in Chapter~\ref{ch:cy-to-chiral}, Theorem~\ref{thm:c3-functor-chain}; the finite Rees hCS--Hall comparison is constructed on $\mathbb{C}^3$ charts and finite DWR/Ran critical atlases by Chapter~\ref{ch:cy3-chain-level-bridge}, while compact Hall doubles require the extra completion, realisation, compact-support, and pairing data isolated below. At $d = 3$ the identification $\kappa_{\mathrm{ch}} = \kappa_{\mathrm{cat}}$ is false for every strict CY$_3$ with $h^{1,0} = 0$: the RHS vanishes by Serre, while the LHS is generically nonzero. The dimension-stratified replacement formula is Conjecture~\textup{\ref{conj:cy-kappa-identification}}, with concrete values catalogued in Theorem~\textup{\ref{thm:kappa-stratification-by-d}} (e.g.\ $\kappa_{\mathrm{ch}}^{\mathrm{Heis}}(\text{quintic}) = -25/3$, $\kappa_{\mathrm{ch}}^{\mathrm{Heis}}(K3 \times E) = 3$, $\kappa_{\mathrm{ch}}^{\mathrm{Heis}}(\text{local }\mathbb{P}^2) = 3/2$; cf.\ Remark~\ref{rem:beauville-kappa-formula-subscript-split}).
\end{remark}

\begin{remark}[Fano obstruction to absolute HPD on \texorpdfstring{$\mathrm{K3}\times E$}{K3x E}; relative HPD over \texorpdfstring{$E$}{E} as fourth Hochschild-pairing path]
\label{rem:cy-cat-hpd-k3e-relative}
Kuznetsov homological projective duality (HPD; Kuznetsov~$2007$~Publ.\ IHES~$105$~Thm.~$1.1$) attaches a dual Lefschetz pair $(Y^\vee,\mathcal A^\vee_\bullet)$ to a \emph{Fano} variety $Y\hookrightarrow\mathbb P^N$ with Lefschetz decomposition $D^b(\Coh(Y))=\langle\mathcal A_0,\mathcal A_1(1),\ldots,\mathcal A_{m-1}(m-1)\rangle$, identifying Kuznetsov components across linear sections. The Fano requirement is load-bearing: the Lefschetz twist is $\mathcal O_Y(1)=\omega_Y^{-1/r}$ for some $r\ge 1$, and this twist is trivial on any strict CY. For $Y=\mathrm{K3}\times E$ the canonical bundle is $\omega_Y\cong\omega_{\mathrm{K3}}\boxtimes\omega_E\cong\mathcal O_Y$ by CY-$2$ triviality of $\omega_{\mathrm{K3}}$ and CY-$1$ triviality of $\omega_E$: absolute HPD does not apply to the product.

The replacement is the \emph{relative} HPD of Kuznetsov~$2018$~Crelle~$736$~\S~$2$ and Kuznetsov--Perry~$2021$ JEMS~$23$ Thm.~$1.1$, applied fibrewise along the second projection $\pi:Y\to E$. Each fibre is a K3 surface; the Kuznetsov--Markushevich~$2009$~J.\ Reine Angew.\ Math.\ $627$ Thm.~$1.1$ / Addington--Thomas~$2014$~Duke Math.\ J.\ $163$ Thm.~$1.1$ criterion identifies generic K3 of degree~$14$ with the Kuznetsov component of a smooth cubic fourfold $Y_4\subset\mathbb P^5$, globalising over~$E$ to a cubic-fourfold fibration $\widetilde Y\to E$ with $\Kuz(\widetilde Y/E)\simeq D^b(\Coh(Y))$. Relative HPD then transports the top-degree categorical Hochschild class $[\chi_3^{\mathrm{cat}}]\in\mathrm{HH}^3_{\mathrm{cat}}(\Kuz(\widetilde Y^\vee/E))$ to the CY-$3$ volume class on~$Y$; under the framed K3-fibre specialisation $\Phi_3^{(K3,E)}$ this becomes the chiral Hochschild cocycle $[\chi_3]\in\ChirHoch^3_{\mathrm{CY\text{-}3}}(\mathbf H_{\Delta_5})$ of Vol~I~Thm.~\texttt{chapters/theory/hochschild\_cohomology.tex} (Vol~I). The value $\langle[\chi_3],[e_3^Y]\rangle_{\Phi_3^{(K3,E)}}=2\,\mathrm{Vol}(E)\cdot(2\pi\mathrm i)^3$ recovers through the Kuznetsov--Perry transport, supplying a fourth categorically-distinct verification of $[\chi_3]\neq 0$ (Vol~I~Thm.~\texttt{chapters/theory/hochschild\_cohomology.tex} (Vol~I)).

Absolute HPD applies to Fano $Y$ only. Relative HPD over a base~$S$ applies to fibrations $Y\to S$ whose fibres are Fano; the base need not be Fano. For $Y=\mathrm{K3}\times E\to E$, neither side is Fano, yet the Kuznetsov--Markushevich upgrade of the K3 fibres to cubic-fourfold Kuznetsov components makes the fibres Fano. Conflating absolute HPD on the CY product with fibrewise absolute HPD on K3 (after the Kuznetsov--Markushevich upgrade) changes the theorem; the relative framework keeps the two scopes separate.
\end{remark}

\providecommand{\Kuz}{\mathrm{Ku}}

\begin{remark}[Drinfeld centre as right adjoint]
\label{rem:cy-cat-drinfeld-center-not-averaging}
The downstream correspondence $\Phi_d$ at $d \geq 3$ produces an $\Eone$-chiral algebra on constructed specialisation loci; the braided $\Etwo$-structure is recovered on the Drinfeld centre $\cZ(\Rep^{\Eone}(\Phi_3^{(\Sigma_2,C)}(\cC)))$ (property~\textup{(U3)} of Theorem~\textup{\ref{thm:phi-platonic}}). The Drinfeld centre is the right adjoint to the forgetful functor $\mathrm{BrMon} \to \mathrm{Mon}$, equivalently the categorified commutant $z(A) = \{a \in A: ab = ba\ \forall b \in A\}$. It differs from a categorified averaging map. The averaging map $\Eone \to \Einf$ loses quantum-group data; the Drinfeld-centre passage $\Eone \to \Etwo$ constructs braiding via half-braidings. This distinction is structural in Vol~III: the chiral-quantum-group content of $\Phi(\cC)$ at $d \geq 3$ is preserved by the centre and lost by symmetric averaging. We therefore reserve ``averaging'' for the $\Eone \to \Einf$ coinvariant projection and ``Drinfeld centre'' for the right-adjoint / half-braiding construction; the two are distinct functors producing distinct output categories.
\end{remark}

\begin{remark}[What \texorpdfstring{$\Phi$}{Phi} does not see, and what CoHA is]
\label{rem:phi-does-not-see}
Not every invariant of $\cC$ survives the passage through $\Phi$. The correspondence retains the cyclic $\Ainf$-structure and the Gerstenhaber bracket; it loses the fine structure of the triangulated / $\Ainf$ category (individual objects, t-structures, stability conditions) visible only via the categorical Cohomological Hall Algebra construction of Chapter~\textup{\ref{ch:toric-coha}}. The Cohomological Hall Algebra $\CoHA(X)$ is the Hochschild cohomology of the quiver-with-potential category, equipped with the Schiffmann--Vasserot--Yang--Zhao multiplication. It is an associative algebra with a Hall product, not a chiral algebra; the identifications ``$\CoHA = \Eone$-chiral'' and ``$\CoHA$ carries a vertex algebra structure'' conflate an associative-algebra output with a chiral-algebra output and so fail as stated. The connection to chiral algebras goes by two separate routes, \emph{not} by a single application of $\Phi$: on one side, the Schiffmann--Vasserot 2013 direct identification $\CoHA(\C^3) = Y^+(\widehat{\mathfrak{gl}}_1)$ (positive half of the affine Yangian of $\widehat{\mathfrak{gl}}_1$) is an equality of associative algebras, not a $\Phi$-application; $\CoHA(\C^3)$ is neither a CY-category nor the $\Phi_3$-input. On the other side, the framed toric specialisation $\Phi_3^{(\Sigma_2,C)}(D^b(\Coh(\C^3)))$ produces a chiral algebra on a reference curve (property~\textup{(U4)} of Theorem~\textup{\ref{thm:phi-platonic}}); specialising via $\SpCh_{\Sigma_2, C}$ with $\Sigma_2 \subset \C^3$ a coordinate plane yields $Y^+(\widehat{\mathfrak{gl}}_1)$ as the Stage-2 shadow on $C$, and the full $\cW_{1+\infty}$ at $c = 1$ then arises from the Drinfeld-centre passage (Remark~\textup{\ref{rem:cy-cat-drinfeld-center-not-averaging}}). The two routes converge on $Y^+(\widehat{\mathfrak{gl}}_1)$ because the Schiffmann--Vasserot CoHA coincides with the $\SpCh$-specialisation of $\PhiFA_3(D^b(\Coh(\C^3)))$ (Kontsevich--Soibelman 2008 \S 4; Schiffmann--Vasserot 2013 \S 6). Writing ``$\Phi(\CoHA(\C^3))$'' would be a type error: $\Phi$ takes CY-categories to chiral algebras and $\CoHA$ is neither.

The surviving numerical data of $\Phi$ are organised by the $\kappa_\bullet$-spectrum (Remark~\textup{\ref{rem:cy-cat-kappa-convention}}): $\kappa_{\mathrm{cat}}$ is the manifold invariant always present, $\kappa_{\mathrm{ch}}$ is the algebraization invariant defined when $\Phi$ exists, $\kappa_{\mathrm{BKM}}$ is the Borcherds-side algebraisation invariant defined for K3-fibered CY$_3$s (Class~A; Proposition~\textup{\ref{prop:bkm-weight-universal}}), and $\kappa_{\mathrm{fiber}}$ is the lattice rank.
\end{remark}

\section{Finite hCS--Hall charts and the compact Hall boundary}
\label{sec:cycat-finite-hcs-hall-boundary}

The CY$_3$ categorical input produces a holomorphic factorisation algebra
before it produces a Hall algebra. A Hall statement enters only after a
critical chart, an orientation line, a Rees/vanishing-cycle realisation,
and the relevant descent topology have been fixed.

\begin{definition}[Finite Rees hCS--Hall comparison chart]
\label{def:cycat-finite-rees-hcs-hall-chart}
Let $X$ be a smooth CY$_3$ and let
\(\sigma\in\mathsf N^{\Ran}_{\leq m}(\mathfrak U;\Gamma_{\leq L})\)
be a finite simplex in a Dolbeault/Weiss/Ran (DWR)-good cover. A
\emph{finite Rees hCS--Hall comparison chart} on $\sigma$ consists of:
\begin{enumerate}[label=\textup{(\roman*)}]
 \item the finite hCS quotient
 \(\Obs_{\hCS}^{q,\leq r,\leq N}(\sigma)_{\leq L}\), with holomorphic
 jet order $\leq N$, renormalisation order $\leq r$, and charge bound
 $\Gamma_{\leq L}$;
 \item a finite cyclic critical model
 \((V_\sigma,W_\sigma,\omega_\sigma)\) over
 \(\Omega^\bullet(\Delta^p)\), obtained by cyclic homological
 perturbation from the hCS BV complex;
 \item the oriented Rees critical Hall complex
 \[
  \Hall_{\sigma}^{\mathrm{Rees},\mathrm{or}}
  =
  \left(
  C^\bullet\!\left(
  \mathfrak g_\sigma,
  \Gamma(V_\sigma,\wedge^\bullet T_{V_\sigma})[[\hbar]]
  \otimes\Omega^\bullet(\Delta^p)
  \right)^{G_\sigma},
  d_{\CE}+\iota_{d_VW_\sigma}
  +\hbar\Delta_{\omega_\sigma}+d_{\Delta^p}
  \right)
  \otimes o_\sigma[s_\sigma](t_\sigma),
 \]
 with determinant-line square root \(o_\sigma\), shift \(s_\sigma\),
 and Tate twist \(t_\sigma\);
 \item a relative degree-zero chain map
 \[
  \theta_\sigma^{\mathrm{rel}}\colon
  B\Obs_{\hCS}^{q,\leq r,\leq N}(\sigma)_{\leq L}
  \longrightarrow
  \Hall_{\sigma}^{\mathrm{Rees},\mathrm{or}}(\Delta^p)
 \]
 compatible with faces, refinements, disjoint union, Thom--Sebastiani,
 and the orientation transports.
\end{enumerate}
This is the finite chart of
Definition~\ref{def:finite-dwr-ran-cyclic-critical-atlas}, rewritten at
the categorical input level. It is a Rees critical complex, not yet the
compact Borel--Moore CoHA unless a monoidal vanishing-cycle realisation
has also been supplied.
\end{definition}

\begin{proposition}[Finite DWR/Ran Rees hCS--Hall comparison]
\label{prop:cycat-finite-dwr-ran-rees-hcs-hall}
Fix finite bounds $(N,r,L,m)$. If a finite DWR/Ran family of charts
\(\{\theta_\sigma^{\mathrm{rel}}\}\) satisfies the compatibility
conditions of Definition~\ref{def:cycat-finite-rees-hcs-hall-chart},
then the sum over simplices
\[
 \Theta_{\hCS\to\Hall}^{\mathrm{Rees},\mathrm{or};N,r,L,m}
 =
 \sum_{0\leq p\leq m}
 \sum_{\sigma\in
 \mathsf N^{\Ran}_{p,\leq m}(\mathfrak U;\Gamma_{\leq L})}
 \int_{\Delta^p}\theta_\sigma^{\mathrm{rel}}
\]
is a Maurer--Cartan element in the finite convolution Lie algebra
\[
 \mathfrak M_{\hCS,\Hall}^{N,r,L,m}
 =
 \Tot_{\sigma}
 \Hom_{\mathrm{cont}}
 \left(
 B\Obs_{\hCS}^{q,\leq r,\leq N}(\sigma)_{\leq L},
 \Hall_{\sigma}^{\mathrm{Rees},\mathrm{or}}
 \right).
\]
Equivalently, it is a continuous multiplicative natural transformation
\[
 \Obs_{\hCS}^{q,\leq r,\leq N}(-)_{\leq L}
 \longrightarrow
 \Hall_{\mathrm{crit}}^{\mathrm{Rees},\mathrm{or},\leq N,\leq L}(-)
\]
on the finite DWR/Ran nerve. The construction covers
\(\mathbb C^3\), finite affine CY$_3$ critical charts, and finite
multi-chart \v{C}ech/Ran diagrams for which the same finite cyclic
critical atlas and orientation data are supplied.
\end{proposition}

\begin{proof}
The cyclic homological perturbation theorem transfers the finite hCS BV
operations to the Hamiltonian \(W_\sigma\). The Fourier--BV
identification of Definition~\ref{def:rees-critical-hall-interlayer}
turns the differential into
\[
 d_{\CE}+\iota_{d_VW_\sigma}+\hbar\Delta_{\omega_\sigma}+d_{\Delta^p}.
\]
The relative master equation is therefore the Maurer--Cartan equation
for \(\theta_\sigma^{\mathrm{rel}}\) in the relative convolution Lie
algebra. Integration over \(\Delta^p\) converts the
\(d_{\Delta^p}\)-term, by Stokes' formula, into the alternating
\v{C}ech/Ran face differential. Summing over all simplices gives
\[
 D_{\mathrm{tot}}\Theta
 +\frac{1}{2}[\Theta,\Theta]=0,
\]
which is exactly the finite descent condition. Multiplicativity is the
Rees Thom--Sebastiani identity for \(W_{\sigma_1}\boxplus W_{\sigma_2}\)
combined with the bar coproduct on hCS observables and the Hall product
on the Rees side. For \(\mathbb C^3\), the chart is the one-vertex
three-loop critical model whose Hall cohomology is the
Schiffmann--Vasserot positive half
\(\CoHA(\mathbb C^3)\cong \Yplus(\widehat{\mathfrak{gl}}_1)\); a finite
affine CY$_3$ critical chart supplies the same data by its
quiver-with-potential critical function; a finite multi-chart diagram
is the preceding construction with the \v{C}ech/Ran face maps retained.
\end{proof}

\begin{definition}[Compact Hall-double realisation data]
\label{def:cycat-compact-hall-double-realisation-data}
For a compact CY$_3$ \(X\), a finite Rees comparison upgrades to a
compact Hall-double comparison only after the following data are fixed:
\begin{enumerate}[label=\textup{(\roman*)}]
 \item a monoidal vanishing-cycle realisation from the Rees critical
 complexes to Borel--Moore critical CoHA, compatible with
 Thom--Sebastiani, proper Hall pushforward, lci pullback, and the
 determinant-line orientation system;
 \item compact-support Beck--Chevalley and DWR/Weiss descent for source
 and target;
 \item a charge/HN and \(\hbar\)-adic completion in which all Hall
 products, coproducts, and inverse limits are continuous;
 \item a Serre-dual negative half, a Cartan half, and a continuous
 Hall--Serre pairing whose radical has been quotiented;
 \item for a Hall--Drinfeld output, centre/half-braiding compatibility
 with the specialised \(\Eone\)-chiral algebra.
\end{enumerate}
A Serre-equivariant quasi-NCCR or a compact critical CoHA character map
is one possible positive-half input to this package. It is not itself a
Hall double.
\end{definition}

\begin{proposition}[Finite positive halves do not define compact Hall doubles]
\label{prop:cycat-compact-hall-double-boundary}
The finite maps of Proposition~\ref{prop:cycat-finite-dwr-ran-rees-hcs-hall}
and the \(K3\times E\) finite Rees maps
\[
 \rho_H^+\colon
 \operatorname{Rees}_{\le H}
 \Yplus_{\mathrm{stalk}}(K3\times E;\mathfrak U)
 \longrightarrow
 \operatorname{Rees}_{\le H}
 \mathcal H_{K3\times E}^{\mathrm{comp},+}
\]
of Proposition~\ref{gluing:prop:compact-coha-rees-maps} construct
positive-half Hall data in finite bounds. They do not construct a
compact Hall--Drinfeld double, a quasi-NCCR realisation map, or a map
from \(\PhiFA_3(\mathrm{Perf}(X))\) to a compact quantum group. Such a
map exists only after the realisation data of
Definition~\ref{def:cycat-compact-hall-double-realisation-data} have
been supplied.
\end{proposition}

\begin{proof}
The Rees comparison of
Proposition~\ref{prop:cycat-finite-dwr-ran-rees-hcs-hall} lands before
vanishing-cycle realisation. Passing to
\(\CoHA_{\mathrm{crit}}^{\mathrm{or}}\) requires the monoidal
realisation functor and the compact-support compatibilities listed in
Definition~\ref{def:cycat-compact-hall-double-realisation-data}. On
\(K3\times E\), Proposition~\ref{gluing:prop:compact-coha-rees-maps}
constructs the compatible finite maps \(\rho_H^+\), but explicitly does
not assert that they are injective, surjective, or primitive-preserving
isomorphisms. Remark~\ref{gluing:rem:no-double-through-hocolim} states
the obstruction in categorical form: a Drinfeld double cannot be
applied to the positive-half homotopy colimit without a completed
monoidal category, dualisability, a non-degenerate Hopf pairing, the
Borcherds bracket comparison, and centre compatibility. The conditional
double of Theorem~\ref{thm:quantum-group-as-positive-geometry-double}
has exactly these hypotheses. Therefore finite DWR/Ran Rees
construction proves the finite positive-half comparison and stops
strictly before the compact Hall-double map.
\end{proof}

\begin{remark}[Cross-volume conventions]
\label{rem:cy-cat-cross-vol-conventions}
Vol~III uses lambda-bracket conventions for the Lie conformal algebras produced by $\Phi$: an OPE of the form $T(z) T(w) \sim (c/2)(z - w)^{-4} + \cdots$ is rewritten as $\{T_\lambda T\} = (c/12) \lambda^3 + \cdots$, absorbing the combinatorial factor $3! = 6$ from the divided-power $\lambda^{(3)} = \lambda^3/3!$. Vol~I uses OPE modes directly; care is required when transporting formulas between volumes (see the concordance). The Hochschild / cyclic invariants of this chapter are convention-independent: they depend only on the chain-level $\Ainf$-structure.
\end{remark}

\subsection{Chain-level \texorpdfstring{$E_d$}{E-d} operad: Fulton--MacPherson configuration spaces}
\label{subsec:cy-cat-fm-config-spaces}

Stage~$1$ of Theorem~\ref{thm:phi-two-stage-factorisation} produces an
$E_d$-holomorphic factorisation algebra $\PhiFA_d(\cC)$. The $E_d$-operad
that appears on the output side has concrete chain-level content: its
arity-$k$ space of operations is the Fulton--MacPherson compactification
$\overline{F}(k, \mathbb{R}^d)^{\mathrm{fm}}$ of the ordered configuration space
of $k$ points in $\mathbb{R}^d$. This subsection records the
Fulton--MacPherson model in the precise form consumed by Kontsevich--Tamarkin
formality and Costello--Gwilliam--Li locality, for both $d = 2$ (the
proved $\PhiFA_2$) and $d = 3$ (the Stage-$1$ half of the CY-A$_3$
construction).

\providecommand{\FMbar}{\overline{F}}
\providecommand{\FMop}{\mathcal{FM}}
\providecommand{\Gerst}{\mathrm{Gerst}}

\begin{definition}[Ordered configuration space; Fulton--MacPherson compactification]
\label{def:cy-cat-fm-config}
For a smooth manifold $M$ of dimension $n$ and $k \geq 1$, the
\emph{ordered configuration space} is
$F(k, M) = \{(x_1, \ldots, x_k) \in M^k : x_i \neq x_j\}$, an open
submanifold of $M^k$ carrying a free $S_k$-action. The
\emph{Fulton--MacPherson compactification}
$\FMbar(k, M)^{\mathrm{fm}}$ (Fulton--MacPherson, \emph{A compactification of configuration spaces}, Ann.\ of Math.\ \textbf{139} (1994) \S 1) is the closure of
$F(k, M)$ in $M^k \times \prod_{|I| \geq 2} \mathrm{Bl}_{\Delta_I}(M^I)$
under the simultaneous inclusions, where $\Delta_I \subset M^I$ is the
thin diagonal indexed by $I \subseteq \{1, \ldots, k\}$ and
$\mathrm{Bl}$ is iterated real blow-up. The space $\FMbar(k, M)^{\mathrm{fm}}$
is a smooth manifold with corners, compact when $M$ is compact, with
$F(k, M)$ embedded as an open dense submanifold.
\end{definition}

\begin{theorem}[Boundary stratification by rooted trees]
\label{thm:cy-cat-fm-strata}
The boundary $\partial \FMbar(k, M)^{\mathrm{fm}}$ stratifies by rooted trees
$T$ with $k$ leaves and no bivalent internal vertices. The stratum
$S_T$ is canonically
\[
 S_T \;\simeq\; F(|v_T|, M) \,\times\, \prod_{v \in \mathrm{int}(T)} F(|v|, \mathbb{R}^n),
\]
where $v_T$ is the root, $|v|$ the number of children at an internal
vertex, and $\mathrm{int}(T)$ the non-root internal vertices. The
$F(|v|, \mathbb{R}^n)$-factors record the relative scale-and-position
data of points colliding into the cluster at $v$, modulo
$\mathbb{R}^n \rtimes \mathbb{R}_{>0}$. For $M = \mathbb{R}^n$ one
passes to the reduced compactification
$\widetilde{\FMbar}(k, \mathbb{R}^n)^{\mathrm{fm}} =
\FMbar(k, \mathbb{R}^n)^{\mathrm{fm}}/(\mathbb{R}^n \rtimes \mathbb{R}_{>0})$,
a smooth manifold with corners of real dimension $nk - n - 1$.
\end{theorem}

\begin{theorem}[Fulton--MacPherson operad; Salvatore--Sinha comparison with little discs]
\label{thm:cy-cat-fm-operad}
The collection
$\FMop_n = \{\widetilde{\FMbar}(k, \mathbb{R}^n)^{\mathrm{fm}}\}_{k \geq 1}$
carries an operad structure in smooth manifolds with corners: the
composition
$\gamma \colon \widetilde{\FMbar}(k, \mathbb{R}^n)^{\mathrm{fm}} \times
 \prod_i \widetilde{\FMbar}(\ell_i, \mathbb{R}^n)^{\mathrm{fm}}
 \to \widetilde{\FMbar}(\sum_i \ell_i, \mathbb{R}^n)^{\mathrm{fm}}$
is a smooth diffeomorphism onto the codimension-$(k-1)$ boundary
stratum $S_{T_{k, \ell_1, \ldots, \ell_k}}$ indexed by the rooted
$2$-level tree with root arity $k$ and children arities $\ell_i$.
The resulting operad is weakly equivalent to the little $n$-discs
operad $\Ddisk_n$ and serves as a $\Sigma$-cofibrant (hence homotopically
good) model for $E_n$; it is weakly equivalent via the
Boardman--Vogt $W$-construction
(Boardman--Vogt, \emph{Homotopy invariant algebraic structures on topological spaces}, Lecture Notes in Math.\ \textbf{347}, 1973; Berger--Moerdijk, \emph{The Boardman--Vogt resolution of operads in monoidal model categories}, Topology \textbf{45} (2006) 807--849, Thm.\ 5.1) to
$W \Ddisk_n$, with length-decorated rooted trees mapping to boundary
strata of $\FMbar(k, \mathbb{R}^n)^{\mathrm{fm}}$.
\end{theorem}

\begin{example}[Cohomology as Gerstenhaber$_n$]
\label{ex:cy-cat-fm-cohomology}
$H^\bullet(F(k, \mathbb{R}^n); \mathbb{Q})$ is generated by the
Arnold classes $g_{ij}$ of degree $n - 1$, modulo $g_{ij}^2 = 0$
and the Arnold triple identity
$g_{ij} g_{jk} + g_{jk} g_{ki} + g_{ki} g_{ij} = 0$
(Arnold, \emph{Mat.\ Zametki} \textbf{5} (1969) 227--231). For $n \geq 2$ this gives
$H^\bullet(\FMop_n; \mathbb{Q}) \simeq \Gerst_n$, the operad of
$(n - 1)$-shifted Poisson algebras: a degree-$0$ graded commutative
product and a Lie bracket $\{-, -\}$ of degree $-(n - 1)$ satisfying
Leibniz (F.~Cohen, \emph{The homology of iterated loop spaces}, Lecture Notes in Math.\ \textbf{533}, 1976; Getzler--Jones, \emph{Operads, homotopy algebra and iterated integrals for double loop spaces}, arXiv:hep-th/9403055). At $n = 3$ the bracket
has degree $-2$, recovering the CY-$3$ shifted-Poisson bracket on
$\HH^\bullet(\cC)$ of Proposition~\ref{prop:cy-bracket-spectrum}.
\end{example}

\begin{theorem}[Kontsevich--Tamarkin $E_n$-formality with $\GRT_1$-torsor]
\label{thm:cy-cat-kt-formality}
For $n \geq 2$ there is a quasi-isomorphism of dg-operads over
$\mathbb{Q}$
\[
 \mathcal{F} \colon
 \Gerst_n \;\xrightarrow{\ \sim\ }\;
 C^{\mathrm{sa}}_\bullet\!\bigl(\FMop_n;\, \mathbb{Q}\bigr),
\]
where $C^{\mathrm{sa}}_\bullet(-; \mathbb{Q})$ is the semi-algebraic
chain complex of Kontsevich--Soibelman (strictly operadic, strictly
multiplicative under Cartesian product). The space of such
$\mathcal{F}$, up to homotopy, is a torsor under the pro-unipotent
Grothendieck--Teichm\"uller group $\GRT_1(\mathbb{Q})$:
\[
 \{\Gerst_n\text{-formality morphisms}\}/{\simeq}
 \;\text{is a }\GRT_1(\mathbb{Q})\text{-torsor}.
\]
The $\GRT_1$-action is realised through Kontsevich graph-integral
weights $W_\Gamma$ evaluated on $\widetilde{\FMbar}(k, \mathbb{R}^n)^{\mathrm{fm}}$
\cite{Kontsevich2003} \S 6.4. Willwacher's theorem
\cite{Willwacher2015} Thm.\ 1.1 identifies
$H^0(\GC_n)$ of the Kontsevich graph complex with the
Grothendieck--Teichm\"uller Lie algebra $\mathfrak{grt}_1$, uniformly
for $n \geq 2$.
\end{theorem}

\begin{construction}[Stage-$1$ $\PhiFA_d(\cC)$ as $\FMop_d$-colimit of Hochschild chains]
\label{constr:cy-cat-phifa-as-fm-colimit}
Fix a smooth complex CY$_d$ variety $X$ and a smooth proper cyclic
$\Ainf$ category $\cC$ of dimension $d$. A \emph{holomorphic $d$-disc}
in $X$ is an open embedding $\varphi \colon B^{2d}_\epsilon
\hookrightarrow X$ from an open polydisc in $\mathbb{C}^d$. Stage-$1$
$\PhiFA_d(\cC)$ assigns to a finite collection of pairwise disjoint
holomorphic $d$-discs $B_1, \ldots, B_k \subset X$ the chain complex
\[
 \PhiFA_d(\cC)(B_1 \sqcup \cdots \sqcup B_k)
 \;\simeq\;
 C^{\mathrm{sa}}_\bullet\!\left(\widetilde{\FMbar}(k, \mathbb{R}^d)^{\mathrm{fm}}; \mathbb{Q}\right)
 \otimes_{C^{\mathrm{sa}}_\bullet(\FMop_d)}
 \HH_\bullet(\cC)^{\otimes k},
\]
the homotopy coend over the $\FMop_d$-operad, with $\HH_\bullet(\cC)$
supplied with its $E_d$-algebra structure through the Kontsevich--Tamarkin
formality morphism of Theorem~\ref{thm:cy-cat-kt-formality}. The
holomorphic refinement is the Dolbeault acyclification over $X$: the
assignment is a sheaf of chain complexes with $\bar\partial$-acyclic
extensions, rigidified by the holomorphic volume $\Omega_X \in
H^0(X, K_X)$ through its role as the BV-rigidifier of the
Costello--Li quantisation scheme
\cite{CostelloLi2016,CostelloLi2020}.
Functorially, the arity-$k$ operation assembles the Hochschild chains
indexed by the external leaves of the rooted trees of
Theorem~\ref{thm:cy-cat-fm-strata} along the operadic composition
$\gamma$.
\end{construction}

\begin{proposition}[Chain-level Dunn additivity $\FMop_3 \simeq \FMop_2 \otimes \FMop_1$]
\label{prop:cy-cat-fm-dunn}
For $m, n \geq 1$ there is a quasi-isomorphism of dg-operads over
$\mathbb{Q}$
\[
 C^{\mathrm{sa}}_\bullet(\FMop_{m+n}; \mathbb{Q})
 \;\xrightarrow{\ \sim\ }\;
 C^{\mathrm{sa}}_\bullet(\FMop_m; \mathbb{Q})
 \otimes_H
 C^{\mathrm{sa}}_\bullet(\FMop_n; \mathbb{Q}),
\]
induced by the coordinate projection
$\widetilde{\FMbar}(k, \mathbb{R}^{m+n})^{\mathrm{fm}} \to
 \widetilde{\FMbar}(k, \mathbb{R}^m)^{\mathrm{fm}} \times
 \widetilde{\FMbar}(k, \mathbb{R}^n)^{\mathrm{fm}}$
(Hadamard tensor of operads). At $(m, n) = (2, 1)$ this recovers
$\FMop_3 \simeq \FMop_2 \otimes \FMop_1$, the chain-level source of the
derived-framing obstruction target
$\HH^{-2}_{E_1}(\HH_\bullet(\cC), \HH_\bullet(\cC))$. Its vanishing is
not a formal consequence of Dunn additivity; Theorem~\ref{thm:derived-framing-obstruction}
requires the comparison and filtration hypotheses stated there. The weak equivalence
is chain-homotopical (a quasi-isomorphism of dg-operads), not a
diffeomorphism of operads in smooth manifolds with corners:
$\FMbar(k, \mathbb{R}^3)^{\mathrm{fm}}$ is not smoothly a product of
$\FMbar(k, \mathbb{R}^2)^{\mathrm{fm}} \times \FMbar(k, \mathbb{R}^1)^{\mathrm{fm}}$.
\end{proposition}

\begin{proposition}[Chain-level Stage-$2$: $\FMop_3 \to \FMop_1$ push-forward along $\Sigma_2$]
\label{prop:cy-cat-fm3-to-fm1}
For an admissible specialisation datum $(\Sigma_2, C)$ with
$\Sigma_2 \subset X$ a smooth real $2$-cycle and $C \subset
X \setminus \Sigma_2$ a holomorphic reference curve, the Stage-$2$
specialisation functor $\SpCh_{\Sigma_2, C}$ is implemented at the
chain level by the operadic push-forward along the fibration
\[
 \widetilde{\FMbar}(k, \mathbb{R}^3)^{\mathrm{fm}}
 \;\longrightarrow\;
 \widetilde{\FMbar}(k, \mathbb{R}^1)^{\mathrm{fm}},
 \qquad
 (x_i^1, x_i^2, x_i^3)_{i=1}^{k}
 \;\mapsto\; (x_i^1)_{i=1}^{k},
\]
projecting each configuration point onto its reference-curve coordinate.
The fibre over a configuration $(x_1^1, \ldots, x_k^1) \in
\widetilde{\FMbar}(k, \mathbb{R}^1)^{\mathrm{fm}}$ is the $\Sigma_2$-direction
compactification of the configuration $(x_i^2, x_i^3)_{i=1}^{k}$,
topologically a $2k$-dimensional fibre. Factorisation homology is
integration along the fibres; the output is an $\Eone$-chiral algebra
on $C$ with operadic operations indexed by $\FMop_1$ composed with
holomorphic pullback. The $\pi_1(\Conf_2(\mathbb{R}^3)) = 0$ symmetry
of the $E_3$-level
(Fadell--Neuwirth, \emph{Configuration spaces}, Math.\ Scand.\ \textbf{10} (1962) 111--118) collapses in this
residual: the non-symmetric quantum-group $R$-matrix datum is
recovered from the Drinfeld centre $\cZ(\Rep^{\Eone}(A_\cC))$
(Remark~\ref{rem:cy-cat-drinfeld-center-not-averaging}), not from
chain-level $\FMop_3$-data.
\end{proposition}

\begin{remark}[Two Fulton--MacPherson spaces: topological $\FMop_3$ vs chiral $\FM_k(\mathbb{C})$]
\label{rem:cy-cat-two-fm-meanings}
The notation $\mathrm{FM}_k(C)$ appears in
Remark~\ref{rem:cy-cat-three-hochschild} for the chiral
Fulton--MacPherson compactification of $k$ points on the curve
$C$: the operadic base space on which the chiral Gerstenhaber bracket
lives through OPE residues. This is distinct from the
$\FMbar(k, \mathbb{R}^3)^{\mathrm{fm}}$ of the present subsection, which
lives on the CY$_3$-ambient side. The relation is
Proposition~\ref{prop:cy-cat-fm3-to-fm1}: the chiral $\FM_k(\mathbb{C})$ is
the Stage-$2$ residual of the ambient $\FMbar(k, \mathbb{R}^3)^{\mathrm{fm}}$
after factorisation-homology integration over $\Sigma_2$ and
holomorphic pullback to $C$. The two Fulton--MacPherson spaces are
\emph{different geometric objects} (one on the ambient manifold,
one on the reference curve) linked by the specialisation functor
$\SpCh_{\Sigma_2, C}$. Swapping the two would also swap the topological Gerstenhaber bracket of Remark~\ref{rem:cy-cat-three-hochschild} with the chiral Gerstenhaber bracket from OPE residues on $\FM_k(\mathbb{C})$, producing an algebra of the wrong type.
\end{remark}

\begin{remark}[$\GRT_1$-torsor and Stage-$1$ canonicity]
\label{rem:cy-cat-grt1-torsor-stage1}
Theorem~\ref{thm:cy-cat-kt-formality} is the chain-level incarnation
of the Stage-$1$ torsor-pointed clause of
Theorem~\ref{thm:phi-two-stage-factorisation}: different choices in
the $\GRT_1(\mathbb{Q})$-torsor produce Stage-$1$ outputs that are
$E_n$-quasi-isomorphic as factorisation algebras, the space of
$E_n$-formality morphisms from $\Gerst_n$ to
$C^{\mathrm{sa}}_\bullet(\FMop_n; \mathbb{Q})$ is a $\GRT_1$-torsor,
and the forgetful map to the space of $E_n$-quasi-isomorphism
classes contracts each torsor orbit to a single point. The
$\GRT_1$-action on graph-integral weights $W_\Gamma$ of
\cite{Kontsevich2003} (\S 6.4) twists the chain-level formality morphism but
preserves the homotopy category of $E_n$-algebras on the nose. In
the holomorphic refinement of
Construction~\ref{constr:cy-cat-phifa-as-fm-colimit} the
$\GRT_1$-action is absorbed into the Costello--Li scheme-dependence
of the BV counter-terms, with fixed observables
\cite{CostelloLi2020}. The torsor is the chain-level incarnation of
the motivic fundamental-group / graph-cohomology structure of
Remark~\ref{rem:cycat-ce-motivic-fundamental-group}.
\end{remark}

\subsection{Fourier--Mukai intertwining of \texorpdfstring{$\PhiFA_3$}{Phi-FA-3} under crepant birational CY\texorpdfstring{$_3$}{-3}}
\label{subsec:cycat-crepant-fm-intertwining}

Two smooth projective CY$_3$ varieties $X$ and $Y$ related by a
crepant birational equivalence $g \colon X \dashrightarrow Y$ have
derived categories connected by a Fourier--Mukai equivalence
$\Phi_g \colon D^b(\Coh(X)) \xrightarrow{\sim} D^b(\Coh(Y))$ with an
explicit kernel $\cK_g \in D^b(\Coh(X \times Y))$. The question
residual to the Stage-$1$ construction of
Theorem~\ref{thm:phi-two-stage-factorisation} is whether this
derived equivalence lifts through $\PhiFA_3$ to a chain-level
quasi-isomorphism of Stage-$1$ factorisation algebras; that is,
whether $\PhiFA_3$ sees the birational equivalence on the CY$_3$ side.

\providecommand{\KM}{\mathrm{KM}}

\begin{theorem}[Stage-$1$ crepant FM-intertwining]
\label{thm:cycat-stage1-crepant-fm-intertwining}
Let $g \colon X \dashrightarrow Y$ be a crepant birational equivalence
of smooth projective CY$_3$ varieties, with Fourier--Mukai kernel
$\cK_g \in D^b(\Coh(X \times Y))$ inducing the derived equivalence
$\Phi_g \colon D^b(\Coh(X)) \xrightarrow{\sim} D^b(\Coh(Y))$. Then
the Stage-$1$ factorisation-algebra assemblies are canonically
$E_3$-quasi-isomorphic:
\[
 \PhiFA_3(\Phi_g) \colon
 \PhiFA_3(D^b(\Coh(X)))
 \;\xrightarrow{\ \sim\ }\;
 \PhiFA_3(D^b(\Coh(Y))),
\]
with intertwining determined by the kernel $\cK_g$ through the
Caldararu Hochschild trace $\Tr_{\cK_g} \colon \HH^\bullet(X) \to
\HH^\bullet(Y)$ (Caldararu~\cite{Caldararu2005}~\S$8$) promoted to an
$E_3$-morphism via the Kontsevich--Tamarkin formality of
Theorem~\ref{thm:cy-cat-kt-formality}. Three input data:
\begin{enumerate}[label=(\roman*)]
 \item \emph{Categorical piece}: the Gerstenhaber-algebra
 quasi-isomorphism $\HH^\bullet(\Phi_g)$ on Hochschild cohomology is
 furnished by the kernel $\cK_g$; for crepant factorisations
 $g = g_n \circ \cdots \circ g_1$ into terminal flops, each
 $\HH^\bullet(\Phi_{g_i})$ is given by Bridgeland/Li, and the
 composition is associative.
 \item \emph{$E_3$-lifting}: the obstruction to lifting
 $\HH^\bullet(\Phi_g)$ to an $E_3$-morphism is in
 $\Ext^2_{\HH^\bullet}(\Gerst_3, \Gerst_3)$, which vanishes over
 $\QQ$ by Fresse~$2011$ \emph{AMS Math.\ Surveys} $217$ Thm~$2$;
 the $\GRT_1$-torsor accumulation across a flop chain is
 cohomologically trivial in the degrees seeing the $E_3$-structure
 (Willwacher~$2015$: $\mathfrak{grt}_1 = H^0(\mathrm{GC}_3)$ is
 generated in odd degree $\geq 3$).
 \item \emph{Holomorphic-volume transport}: the phase
 $e^{i\theta_g} \in \mathrm{U}(1)$ relating $g^*\Omega_Y$ to
 $\Omega_X$ is trivial because $g$ factors through connected
 components of the birational-automorphism group
 (Matsuki~$2002$ Ch.~$1$); hence $g^*\Omega_Y = \Omega_X$ strictly,
 and the Costello--Li one-loop BCOV curving
 $\alpha_{\mathrm{BCOV}} = (\chi_{\mathrm{top}}/24)[\Omega_X]^{0,1}$
 transports rigidly (Batyrev~$1999$: $\chi_{\mathrm{top}}$ is a
 birational invariant of smooth projective CY$_d$).
\end{enumerate}
\end{theorem}

\begin{proof}[Proof sketch]
For $g$ a simple flop, Bridgeland~$2002$ gives an explicit kernel
$\cK_g = \cO_{X \times_W X^+}$ on the fibre product over the
flopping contraction $W$, and the induced Hochschild
quasi-isomorphism $\HH^\bullet(\Phi_g)$ preserves the
Gerstenhaber bracket (Bodzenta--Bondal~$2015$). For $g$ a composition
$g = g_n \circ \cdots \circ g_1$ of terminal flops
(Kawamata--Matsuki~$2008$ Thm~$1$), the intertwining
$\PhiFA_3(\Phi_g) = \PhiFA_3(\Phi_{g_n}) \circ \cdots \circ
\PhiFA_3(\Phi_{g_1})$ assembles by associativity of factorisation-algebra
morphisms (Costello--Gwilliam~\cite{CostelloGwilliam} Prop.~$7.4$).
The $E_3$-lift from Gerstenhaber quasi-isomorphism to
$E_3$-quasi-isomorphism has vanishing obstruction (Fresse loc.\ cit.),
and the cumulative $\GRT_1$-twist is in the commutator subalgebra of
$\mathfrak{grt}_1$, which has no generators in cohomological degrees
seeing the $E_3$-structure (Willwacher~$2015$
\emph{Invent.}~$200$ Thm~$1.1$). The holomorphic-volume phase is
$+1$ because $\mathrm{Bir}(X, Y)/\mathrm{Iso}(X, Y)$ is connected
(Matsuki loc.\ cit.) and characters of connected groups are trivial.
$\chi_{\mathrm{top}}$-birational-invariance gives rigid transport of
$\alpha_{\mathrm{BCOV}}$.
\end{proof}

\begin{remark}[Chain-level transport via the common resolution]
\label{rem:cycat-atiyah-common-resolution}
For the Atiyah flop $g \colon X \dashrightarrow X^+$ (replacing a
$(-1,-1)$-curve $\mathbb{P}^1 \subset X$ with its dual), the common
resolution $Z \xrightarrow{p_X} X$, $Z \xrightarrow{p_Y} X^+$ is the
blow-up of the conifold point. The Stage-$1$ intertwining is realised
at the chain level by pull-push along the common resolution:
\[
 \PhiFA_3(\cC_{X^+}) \;\simeq\; (p_Y)_* (p_X)^* \, \PhiFA_3(\cC_X),
\]
with the Atiyah-flop kernel $\cO_{X \times_W X^+}$ corresponding to a
class in $\HH^\bullet(\cC_Z)$ supported on the exceptional divisor.
The pagoda flop (iterated Atiyah flop) closes by associativity of the
kernel convolution
$\cK_g = \cK_{g_n} \star \cdots \star \cK_{g_1}$ on the fibre product
(Bondal--Orlov~$1995$~\S $3$).
\end{remark}

\begin{remark}[Abuaf non-flop crepant equivalence: direct kernel computation]
\label{rem:cycat-abuaf-non-flop}
Abuaf~\texttt{arXiv:1510.05257} constructs a derived equivalence
$\Phi_{\mathrm{Abuaf}} \colon D^b(\widetilde Z) \simeq D^b(W)$ between
a small resolution of a weighted-projective hypersurface $\widetilde Z$
and a smooth projective variety $W$ obtained by a Mukai flop on a
Grassmannian-bundle contraction. The two spaces $\widetilde Z$ and $W$
are \emph{not} related by a chain of Kawamata--Matsuki flops: they have
distinct birational models at the generic point of the exceptional
divisor. The Kawamata--Matsuki factorisation strategy of
Theorem~\ref{thm:cycat-stage1-crepant-fm-intertwining} therefore does
not apply directly. The $\PhiFA_3$-intertwining nonetheless holds by
direct kernel computation: the Abuaf kernel $\cK_{\mathrm{Abuaf}}$ sits
in the same $K$-theoretic class as the convolution of a generalised
Mukai-flop pair on the auxiliary Grassmannian-bundle structure, and
the induced Caldararu trace $\Tr_{\cK_{\mathrm{Abuaf}}}$ is a
Gerstenhaber quasi-isomorphism promoted to an $E_3$-quasi-isomorphism
by the same Fresse vanishing. The holomorphic-volume phase is $+1$
(Abuaf loc.\ cit.\ Thm~$2$: the Mukai reflection on the Grassmannian
fibre is orientation-preserving). Hence crepant equivalence in the
Abuaf non-flop example likewise induces a Stage-$1$
$E_3$-quasi-isomorphism of $\PhiFA_3$.
\end{remark}

\begin{remark}[Four-\texorpdfstring{$\kappa_\bullet$}{kappa-bullet} discipline under crepant birational maps]
\label{rem:cycat-four-kappa-crepant}
Under crepant birational $g \colon X \dashrightarrow Y$:
$\kappa_{\mathrm{cat}}(X) = \chi(\cO_X) = \chi(\cO_Y) = \kappa_{\mathrm{cat}}(Y)$
(Kollar~$1986$: crepant preserves $\chi(\cO)$); the Stage-$1$
$E_3$-quasi-isomorphism of
Theorem~\ref{thm:cycat-stage1-crepant-fm-intertwining} implies
$\kappa_{\mathrm{ch}}(\Phi_3(\cC_X)) = \kappa_{\mathrm{ch}}(\Phi_3(\cC_Y))$
after fixing a common admissible Stage-$2$ specialisation $(\Sigma_2, C)$
compatible with the birational map (H1--H4 of
Theorem~\ref{thm:cy-to-chiral-d3}). $\kappa_{\mathrm{fiber}}$ is
\emph{not} birational-invariant in general: the Atiyah flop changes
the divisor-class lattice on the exceptional locus, hence changes
$\kappa_{\mathrm{fiber}}$. $\kappa_{\mathrm{BKM}}$ is a Stage-$2$
automorphic invariant and is not defined at the Stage-$1$ level;
its transport under $g$ requires a compatible Stage-$2$ framing.
\end{remark}

\subsection{Arity-\texorpdfstring{$\geq 4$}{>= 4} Poincar\'e self-pairing on \texorpdfstring{$\FMbar(k, \mathbb{R}^3)^{\mathrm{fm}}$}{FM-bar(k, R^3)}}
\label{subsec:cycat-arity-ge-4-poincare}

The Ayala--Francis operadic Koszul self-duality
$E_3^! \simeq E_3\{3\}$ (Ayala--Francis~$2014$ \emph{arXiv:$1409.2478$}
Thm~$2.1$) is witnessed chain-level by the Poincar\'e pairing on the
Fulton--MacPherson compactification $\widetilde{\FMbar}(k, \mathbb{R}^3)^{\mathrm{fm}}$,
a smooth manifold-with-corners of real dimension $3k - 3$. At arity
$k = 2$ the reduced space is $S^2$ (degree $2 = n - 1$ Arnold class,
top degree $2$ on $S^2$); at $k = 3$ the reduced space has real
dimension $6$ and the Poincar\'e pairing gives the Arnold-triple
product of degree $6$. For $k \geq 4$ the configuration space carries
a richer combinatorial structure; the closed-form primitive class and
its self-pairing are the subject of this subsection.

\begin{theorem}[Primitive self-dual class and Poincar\'e self-pairing at arity $k$]
\label{thm:cycat-arity-k-primitive-class}
For $k \geq 2$, the reduced Fulton--MacPherson compactification
$\widetilde{\FMbar}(k, \mathbb{R}^3)^{\mathrm{fm}}$ carries a unique
$\Sym_k$-equivariant top-Arnold-degree primitive class
\[
 \omega_k^{\mathrm{prim}}
 \;\in\;
 H^{2(k-1)}\!\bigl(F(k, \mathbb{R}^3); \mathbb{Q}\bigr)^{\Sym_k}
\]
whose Poincar\'e self-pairing equals
\[
 \bigl\langle \omega_k^{\mathrm{prim}}, \omega_k^{\mathrm{prim}, \vee}\bigr\rangle_{\FMop_3(k)}
 \;=\; k!.
\]
Explicit values:
$k = 2$: pairing $= 2$; $k = 3$: $= 6$; $k = 4$: $= 24$;
$k = 5$: $= 120$. The primitive class at arity $k$ is the
$\Sym_k$-symmetrised Pfaffian-like combination
\[
 \omega_k^{\mathrm{prim}}
 \;=\;
 \frac{1}{k!}
 \sum_{\sigma \in \Sym_k}
 \mathrm{sgn}(\sigma) \,
 g_{\sigma(1)\sigma(2)} g_{\sigma(3)\sigma(4)}
 \cdots g_{\sigma(2\lfloor k/2 \rfloor - 1)\sigma(2\lfloor k/2 \rfloor)}
 \cdot (\mathrm{Sym\text{-}completion}),
\]
where $g_{ij}$ is the Arnold class of degree $n - 1 = 2$, modulo the
Arnold triple relations
$g_{ij} g_{jk} + g_{jk} g_{ki} + g_{ki} g_{ij} = 0$
(Arnold~$1969$). The $k!$-normalisation is the $\Sym_k$-orbit count on
the ordered configuration space $\Conf_k(\mathbb{R}^3)$.
\end{theorem}

\begin{proof}[Proof]
The ordered configuration space $F(k, \mathbb{R}^3)$ has rational
cohomology given by Cohen's formula
$P_k^{(n)}(t) = \prod_{j=1}^{k-1}(1 + j\, t^{n-1})$
(F.~Cohen~$1976$ \emph{LNM} $533$ Thm~$5.1$). At $n = 3$ and each
$k$, the top-degree component is generated by
the Pfaffian-like primitive
$\omega_k^{\mathrm{prim}}$; this is the unique top-Arnold-degree
$\Sym_k$-equivariant class by the Pfaffian recognition theorem
(Knutson--Miller~$2005$
\emph{Ann.\ Math.}~$161$ \S~$4$, applied to the Arnold algebra).
Poincar\'e duality on $\widetilde{\FMbar}(k, \mathbb{R}^3)^{\mathrm{fm}}$
(oriented smooth manifold-with-corners of real dimension $3(k-1)$;
Salvatore--Sinha~$2001$ \emph{Progr.\ Math.\ $196$}) gives a pairing
$H^{2(k-1)} \otimes H^{(k-1)} \to \mathbb{Q}$, whose
$\Sym_k$-equivariant restriction to the primitive subspace is
the $\Sym_k$-orbit count, which is $k!$ on the ordered
configuration space. For $k = 4$: the lemma follows by direct
arity-four computation.
For general $k$: induction via the Dunn-additivity factorisation
$\FMop_3 \simeq \FMop_2 \otimes \FMop_1 / (\text{braid diagonal})$
(Proposition~\ref{prop:cy-cat-fm-dunn}), with the $\FMop_2$-pairing
at arity $k$ equal to $k!$ (Getzler--Jones~$1994$ eq.~$5.3$) and the
$\FMop_1$-pairing at arity $k$ equal to $k!$ (Kapranov--Manin~$2001$
for real-line configuration spaces), reduced by the braid-diagonal
factor $1/k!$ to recover $k!$ on $\FMop_3$.
\end{proof}

\begin{corollary}[Consistency with operadic Koszul self-duality shift]
\label{cor:cycat-arity-k-koszul-consistency}
The $k!$-normalisation of
Theorem~\ref{thm:cycat-arity-k-primitive-class} is consistent with
the Ayala--Francis operadic Koszul self-duality
$E_3^! \simeq E_3\{3\}$: the manifold-dimension pairing on
$\widetilde{\FMbar}(k, \mathbb{R}^3)^{\mathrm{fm}}$ has degree
$3(k-1)$, which under the Koszul shift $\{3\}$ corresponds to an
operadic pairing of degree $3(k-1) - 3(k-1) = 0$ on the operad and
degree $3$ on the target algebra. At each arity $k$ the shift
degree accounting is consistent: the operadic degree at arity $k$ is
$n(k-1) = 3(k-1)$, matching the manifold dimension of the
compactification, after the uniform shift of $n = 3$.
\end{corollary}

\begin{remark}[Relation to Stage-$1$ $\PhiFA_3$ self-pairing]
\label{rem:cycat-arity-k-phi-fa-pairing}
The Poincar\'e self-pairings of
Theorem~\ref{thm:cycat-arity-k-primitive-class} are the operadic
shadow of the CY$_3$ cyclic self-pairing on
$\PhiFA_3(\cC)$ (Stage-$1$ self-pairing theorem):
the CY$_3$ cyclic pairing of degree $3$ on $\HH^\bullet(\cC)$ lifts
through the Ayala--Francis factorisation-homology map
$\int_{\mathbb{R}^3}(\cA) \simeq \Obs_{E_3}(\cA)$
(Ayala--Francis~$2014$ Thm~$4.2$) to a pairing on the arity-$k$
observables, with the arity-$k$ multiplicity $k!$ exactly tracking
the $\Sym_k$-equivariance of the factorisation-algebra structure.
The consistency with Koszul self-duality
(Corollary~\ref{cor:cycat-arity-k-koszul-consistency}) is what
realises the CY$_3$ cyclic pairing as the Stage-$1$ incarnation of
the operadic shift $E_3^! \simeq E_3\{3\}$.
\end{remark}

\section{CY\texorpdfstring{$_2$}{-2} and CY\texorpdfstring{$_3$}{-3} at K3: two realisations, one chiral bialgebra}
\label{sec:cycat-K3-two-realisations}

\begin{remark}[CY\texorpdfstring{$_2$}{-2} versus CY\texorpdfstring{$_3$}{-3} at the K3 base point]
\label{rem:cycat-K3-two-sides}
Two distinct CY-category realisations meet at K3: the CY-$2$ category
$D^b\mathrm{Coh}(K3)$, Serre functor $S = [2]\otimes(-)$, carrying the
$\widetilde M_{24}$ Schur cocycle of order $6$; and the CY-$3$ category
$\mathrm{CoHA}_{K3\times E}$ of the elliptic-fibred total space, Serre
shift $[3]$, extended by Davison. The chiral bialgebra
$\mathbf H_{\Delta_5}$ bridges both: Etingof--Kazhdan quantisation of
the BKM bialgebra attached to $D^b\mathrm{Coh}(K3)$ on the CY-$2$ side,
Hall structure supplied by the Schiffmann--Vasserot CoHA on $K3\times E$
on the CY-$3$ side. See Chapter~\ref{ch:k3-chiral-bialgebra-platonic};
primary sources Kontsevich--Soibelman 2008, Davison 2017,
Schiffmann--Vasserot 2012.
\end{remark}

\section{The Conway \texorpdfstring{$V^{s\natural}$}{V-snatural}-locus of the CY-category landscape}
\label{sec:cycat-witten-conway}

The $\{\Phi_d\}$ correspondence programme (Chapter~\textup{\ref{ch:cy-to-chiral}})
admits a distinguished \emph{super-CY fibre} at the K3 moduli
$T^4/\mathbb{Z}/2$-orbifold point, where the super-moonshine
module of Duncan is realised as a CY-categorical limit of
the Fukaya category of the K3 at that point.

\begin{remark}[Taormina--Wendland super-Fukaya locus at the K3 \texorpdfstring{$T^4/\mathbb{Z}/2$}{T-4/Z/2}-orbifold]
\label{rem:cycat-taormina-wendland-locus}
\index{Taormina-Wendland!super-Fukaya locus}
\index{Duncan module!CY-category realisation}
Taormina--Wendland 2011 (\emph{Commun.~Number Theory Phys.}~7, 527--587,
arXiv:1107.3834) and Taormina--Wendland 2013 (\emph{J.\ High Energy
Phys.}~04~102, arXiv:1307.3892) identified a distinguished point in
K3 moduli; the \emph{Kummer $T^4/\mathbb{Z}/2$-orbifold}; at
which the K3 sigma-model chiral algebra contains a sub-chiral-algebra
canonically isomorphic to Duncan's Conway super-VOA $V^{s\natural}$
(central charge $c = 12$) up to an $\mathrm{Co}_0$-equivariant
embedding. On the CY-category side, this distinguished locus
corresponds to a Fukaya-categorical relation
\[
 \mathrm{Fuk}(K3)\bigr|_{T^4/\mathbb{Z}/2}
 \;\xhookrightarrow{\;\Psi_{\mathrm{TW}}\;}\;
 \mathrm{Fuk}_{\mathrm{super}}(\Lambda_{24})
 \;\;\supset\;\;
 V^{s\natural},
\]
where $\mathrm{Fuk}_{\mathrm{super}}(\Lambda_{24})$ is the
super-Fukaya category of the Leech torus
$T^{24} = \mathbb{R}^{24}/\Lambda_{24}$ equipped with the Duncan
$N{=}1$-super structure.
The CY-to-chiral functor $\Phi$ sends
$\mathrm{Fuk}(K3)|_{T^4/\mathbb{Z}/2}$
to the K3 chiral bialgebra $\mathbf{H}_{\Delta_5}$ restricted to
the Conway locus (Vol.~III
Theorem~\textup{\ref{thm:bkm-taormina-wendland-conway-embedding}});
simultaneously, $\Phi$ restricted to the super-Fukaya category
of $\Lambda_{24}$ produces the Conway chiral bialgebra
$\mathbf{H}_{\mathrm{Conway}}$ of Vol.~III
Conjecture~\textup{\ref{thm:bkm-conway-psi-fifth-image-resolved}}. The
compatibility
\[
 \Phi\bigl(\mathrm{Fuk}(K3)|_{T^4/\mathbb{Z}/2}\bigr)
 \;\hookrightarrow\;
 \Phi\bigl(\mathrm{Fuk}_{\mathrm{super}}(\Lambda_{24})\bigr)
 \;=\;
 \mathbf{H}_{\mathrm{Conway}}
\]
is \emph{both} a CY-category restriction (along the
$T^4/\mathbb{Z}/2$-orbifold sublocus) \emph{and} a $\Psi$-functorial
embedding (fifth $\Psi$-image restricted to the K3 flagship).
Primary: Taormina--Wendland 2011, 2013;
Duncan 2007 \emph{Math.~Res.\ Lett.}~14;
Vol.~III Theorem~\textup{\ref{thm:bkm-taormina-wendland-conway-embedding}}.
\end{remark}

\begin{remark}[\texorpdfstring{$\kappa_\bullet$}{kappa-.}-spectrum on the Conway row: super-extension discipline]
\label{rem:cycat-kappa-spectrum-conway}
\index{kappa spectrum!Conway row}
\index{super-kappa discipline}
The Conway row of the conjectural five-entry $\Psi$-landscape
(Vol.~III Conjecture~\textup{\ref{thm:bkm-conway-psi-fifth-image-resolved}})
distinguishes itself from the four bosonic rows (K3, Enriques,
Monster, Fake-Monster) at the level of the $\kappa_\bullet$-spectrum
of Remark~\textup{\ref{rem:cy-cat-kappa-convention}}. Writing
$\kappa^{\mathrm{super}}_{\mathrm{ch}}$ for the
super-chiral modular characteristic (supertrace-weighted):
\begin{itemize}
 \item $\kappa_{\mathrm{ch}}(V^{s\natural}) = c/2 = 6$ in the
 \emph{super-normalisation} (per the Vol.~I Theorem~D
 $\kappa_{\mathrm{ch}} = c/2$ on free-field class, adapted to super via
 supertrace on $V^{s\natural}_{1/2} = 0$ with all mass at
 integer weight; cf.~Vol.~I chapter on chiral moonshine unified).
 \item $\kappa_{\mathrm{cat}}$ is \emph{not a CY-categorical
 invariant of $V^{s\natural}$} directly: $V^{s\natural}$ is not the
 $\Phi$-image of a $D^b(\mathrm{Coh}(X))$ or
 $\mathrm{Fuk}(M)$ input for a compact complex / symplectic
 variety. The Taormina--Wendland embedding of
 Remark~\ref{rem:cycat-taormina-wendland-locus} realises
 $V^{s\natural}$ as the $\Psi$-image of a \emph{super}-Fukaya
 category $\mathrm{Fuk}_{\mathrm{super}}(\Lambda_{24})$,
 which sits outside the bosonic $\mathrm{CY}_d$ landscape of
 this chapter. The resulting $\kappa^{\mathrm{super}}_{\mathrm{cat}}$
 is the super-Euler-characteristic
 $\chi_{\mathrm{super}}(\mathcal{O}_{T^{24}/\Lambda_{24}})$ of
 the Leech torus (super-normalisation), requiring an
 extension of the four-entry $\kappa_\bullet$-convention to
 super-categorical inputs.
 \item $\kappa_{\mathrm{BKM}}(\mathbf{H}_{\mathrm{Conway}})$ is
 the Borcherds-weight invariant of the super-automorphic form
 $\Phi_{\mathrm{Conway}} = (\Theta_{\Lambda_{24}}/\eta^{24})^{1/2}$.
 Duncan 2007 Thm~4.2 computes its super-weight $c_1^{\mathrm{super}}(0)/2$
 as the sum over the $24$ Leech simple roots.
\end{itemize}
On the three-faces-identity line of the Conway row, the universal identity $\hbar^2 K = -1$ specialises on the super-extension $V^{s\natural}$ at $c = 12$ with $K_{\mathrm{Conway}} = 2$ and $\hbar^2_{\mathrm{Conway}} = -1/2$ (Duncan 2007 \S 6 diamond). The distinction between the three-faces-identity lattice conductor $K = 2$ (positive-signature count on the super-metaplectic input; bosonic Leech $K = 48$ out of scope) and the Borcherds super-weight value is a $K$-vs-$\kappa_{\mathrm{BKM}}$ separation that the super-Borcherds-product expansion must resolve on its own terms; the super-product of the square-root expression $(\Theta_{\Lambda_{24}}/\eta^{24})^{1/2}$ is not expanded in this chapter.

The coincidence $\kappa^{\mathrm{super}}_{\mathrm{ch}}(V^{s\natural})
= \kappa_{\mathrm{ch}}(V^\natural)/2 = 6$ is the super-extension
face of the Monster / Conway lattice diamond
(Vol.~I Remark~\textup{\texttt{chapters/examples/lattice\_foundations.tex} (Vol~I)});
it is \emph{not} an equality of the same invariant evaluated on
two algebras, but the $\mathbb{Z}/2$-super-graded halving of a
bosonic invariant. Conflating the two violates
the $\kappa_\bullet$-discipline of
Remark~\ref{rem:cy-cat-kappa-convention}: Monster and Conway
inhabit two different rows of the $\Psi$-landscape and
carry two different $\kappa^{(\bullet)}_{\mathrm{ch}}$ invariants
(bosonic vs.\ super), coincidentally equal only on the common
$c = 12$ face. Primary: Duncan 2007 Thm~4.2;
Vol.~I Chapter on chiral moonshine unified
(\texttt{chiral\_moonshine\_unified.tex} \S4 for
$\kappa_{\mathrm{ch}}(V^{s\natural})$);
Vol.~III Conjecture~\textup{\ref{thm:bkm-conway-psi-fifth-image-resolved}};
Scheithauer 2006.
\end{remark}

\section{Hodge structure on Chevalley--Eilenberg cohomology of \texorpdfstring{$\mathfrak{g}_{\Delta_5}$}{g-Delta-5}}
\label{sec:cycat-ce-hodge-delta5}

The BKM Lie superalgebra $\mathfrak{g}_{\Delta_5}$ (Chapter~\textup{\ref{ch:k3-chiral-bialgebra-platonic}})
carries two gradings inherited from K3 geometry: the root-lattice
grading on $\Lambda^{2,1}_{II} \subset \mathrm{II}_{3,19}$
(separating real and imaginary roots) and the Mukai lattice grading
$\mathrm{II}_{4,20}$ (the horizontal 24-direction). Both gradings
have Hodge-theoretic origins traceable to the K3 period map. The
Chevalley--Eilenberg cohomology inherits a \emph{mixed Hodge structure}
whose weight and Hodge filtrations are pulled back from these K3
periods via the automorphic Borcherds lift.

\begin{theorem}[CE cohomology as sum over positive roots]
\label{thm:cycat-ce-sum-positive-roots}
For any Borcherds--Kac--Moody Lie superalgebra $\mathfrak{g}$ with
triangular decomposition $\mathfrak{g} = \mathfrak{n}_- \oplus
\mathfrak{h} \oplus \mathfrak{n}_+$ (super-grading from imaginary
simple roots included), the Chevalley--Eilenberg cohomology of the
positive nilpotent subalgebra decomposes as
\[
 H^k_{\mathrm{CE}}(\mathfrak{n}_+) \;=\;
 \bigoplus_{\substack{\alpha \in \Delta_+^{\mathrm{ord}}\\ \ell(\alpha) = k}}
 \Lambda^{k}\mathfrak{g}_\alpha^*
\]
where $\Delta_+^{\mathrm{ord}}$ is the set of positive roots ordered
by Weyl length $\ell$ (Kostant's cubic Dirac operator). For the full
algebra $\mathfrak{g}_{\Delta_5}$, the denominator identity
\[
 \Delta_5(\rho, \tau, z) \;=\;
 \sum_{w \in W} \varepsilon(w) \cdot e^{w\rho} \cdot
 \prod_{\alpha \in \Delta_+^{\mathrm{im}}}
 (1 - e^{-\alpha})^{\mathrm{mult}(\alpha)}
\]
recovers CE Euler-character identities: taking the super-trace over
all $k$ gives $\chi_{\mathrm{CE}}(\mathfrak{g}_{\Delta_5}) = \Delta_5$
as a Siegel form on $\mathbb{H}_2$. Multiplicities of imaginary simple
roots are $\mathrm{mult}(\alpha) = c_{\mathrm{K3}}(4nm - \ell^2)$ where
$c_{\mathrm{K3}}$ are the Fourier coefficients of the K3 elliptic
genus $\phi_{0,1}^{\mathrm{K3}}$ (Eguchi--Ooguri--Tachikawa 2011).
\end{theorem}

\begin{definition}[Hodge filtration on CE cohomology]
\label{def:cycat-ce-hodge-filtration}
Writing $H^k_{\mathrm{CE}}(\mathfrak{g}_{\Delta_5}) \otimes_\mathbb{Z}
\mathbb{C}$, define the \emph{Hodge filtration} $F^\bullet H^k$ via
the Mukai lattice projection $\pi_{\mathrm{Muk}}: \mathfrak{g}_{\Delta_5}
\to \mathrm{II}_{4,20} \otimes \mathbb{C}$ composed with the K3
period isomorphism
\[
 \mathrm{Per}^{\mathrm{Muk}} \colon
 \mathrm{II}_{4,20} \otimes \mathbb{C}
 \xrightarrow{\ \sim\ }
 H^{2,0}(\mathrm{K3}) \oplus H^{1,1}(\mathrm{K3}) \oplus H^{0,2}(\mathrm{K3})
 \oplus H^{0,0}(\mathrm{K3}) \oplus H^{2,2}(\mathrm{K3}),
\]
where the four-dimensional positive plane $H^{2,0} \oplus H^{0,0} \oplus
H^{2,2}$ carries the holomorphic K3 period $\omega_{K3}$ and its
complex conjugate. Put $F^p H^k_{\mathrm{CE}}$ equal to the image
in $H^k$ of $\bigoplus_{p' \geq p}(H^{p',k-p'})$ under this
period isomorphism, pulled back to each $\mathfrak{g}_\alpha^*$
via the root-lattice embedding $\alpha \hookrightarrow
\mathrm{II}_{4,20}$. Define the \emph{weight filtration} $W_\bullet$ by
$W_k = \bigoplus_{\alpha: \alpha^2 > 0} \Lambda^k \mathfrak{g}_\alpha^*$
for real roots (weight $k$) and
$W_{2k} \supset \bigoplus_{\alpha: \alpha^2 \leq 0, \mathrm{mult}(\alpha)}
\Lambda^k \mathfrak{g}_\alpha^*$ for imaginary roots, with the
mixing on the imaginary cone read off from $\phi_{0,1}^{K3}$ Fourier
coefficients.
\end{definition}

\begin{theorem}[Real roots are pure Tate \texorpdfstring{$(k,k)$}{(k,k)}; imaginary roots are MHS]
\label{thm:cycat-ce-hodge-real-imag}
With the filtrations of Definition~\textup{\ref{def:cycat-ce-hodge-filtration}},
the Hodge structure on each graded piece of CE cohomology splits:
\begin{enumerate}[label=(\roman*)]
 \item For each real root $\alpha \in \Delta_+^{\mathrm{re}}$
 (satisfying $\alpha^2 = 2$, $\mathrm{mult}(\alpha) = 1$), the
 graded piece
 $\mathrm{gr}^k_F H^k_{\mathrm{CE},\alpha}(\mathfrak{g}_{\Delta_5})
 \simeq \mathbb{Q}(-k)$
 is \emph{pure Tate} of weight $2k$, type $(k,k)$. This is forced
 by: real roots sit in the positive-definite $(2,1)$-signature
 hyperbolic Cartan $\Lambda^{2,1}_{II}$, which does not intersect
 the K3 holomorphic period $\omega_{K3} \in H^{2,0}$; their
 contribution to $\Delta_5$ is the Weyl denominator
 $\prod_\alpha (1 - e^{-\alpha})$, a pure algebraic factor.
 \item For each imaginary root $\alpha \in \Delta_+^{\mathrm{im}}$
 (satisfying $\alpha^2 \leq 0$, $\mathrm{mult}(\alpha) =
 c_{\mathrm{K3}}(4nm - \ell^2)$), the graded piece carries a genuinely
 \emph{mixed Hodge structure}: $\mathrm{gr}^W_n H^k_{\mathrm{CE},\alpha}$
 is non-trivial in weights $n \in \{0, 1, \ldots, 2k\}$, with
 weight-0 piece coming from the $H^{1,1}(\mathrm{K3})$ contribution to
 $\phi_{0,1}^{\mathrm{K3}}$ and weight-$2k$ piece from the pairing
 of $H^{2,0} \otimes H^{0,2}$ with $\alpha^{\otimes k}$. Concretely, the
 weight filtration is pulled back from the Mukai lattice via
 \[
 \phi_{0,1}^{\mathrm{K3}}(\tau, z) \;=\;
 \sum_{n \geq -1} c_{\mathrm{K3}}(n) \, \theta_{\mathrm{K3}}^{(n)}(\tau, z),
 \qquad
 W_{2n}/W_{2n-2} \;\simeq\; \theta_{\mathrm{K3}}^{(n)}\text{-component}.
 \]
\end{enumerate}
The total Hodge structure on $H^\bullet_{\mathrm{CE}}(\mathfrak{g}_{\Delta_5})$
is a direct sum of (i) pure Tate-$(k,k)$ pieces for real roots and
(ii) mixed pieces for imaginary roots, with the mixing controlled by
the Hodge decomposition $\phi_{0,1}^{\mathrm{K3}} = \phi^{(2,0)} +
\phi^{(1,1)} + \phi^{(0,2)}$ of the K3 elliptic genus into Hodge
components.
\end{theorem}

\begin{proof}[Proof sketch]
For (i): the real-root reflection hyperplanes are parallel to
$H^{1,1}(\mathrm{K3})_{\mathbb{R}}$ inside the Mukai lattice after
complexification. Under $\mathrm{Per}^{\mathrm{Muk}}$ the reflection
$\alpha \mapsto w(\alpha) = \alpha - \langle\alpha, v\rangle v$ is
a real operator preserving the $(2,20)$-signature hyperplane; its
eigenspaces are one-dimensional and carry no Hodge-filtration
obstruction, giving pure type $(k,k)$ on $k$-fold wedge.
For (ii): the imaginary-root multiplicity $\mathrm{mult}(\alpha) =
c_{\mathrm{K3}}(4nm-\ell^2)$ expands Fourier-theoretically as a sum
over $(n,\ell,m)$-refined K3 Poincar\'e sums, and Bruinier's
singular-theta lift (Bruinier LNM 1780 \S 2.2) decomposes the
multiplicity via the Hodge splitting of $\phi_{0,1}^{\mathrm{K3}}$
into $(H^{2,0}, H^{1,1}, H^{0,2})$-pieces, inducing the weight
filtration $W_\bullet$ on each $\Lambda^k \mathfrak{g}_\alpha^*$.
The splitting is strict because the Borcherds singular-theta lift
is functorial in the Hodge filtration of its input.
\end{proof}

\begin{corollary}[Motivic origin: K3 periods carry the Hodge data]
\label{cor:cycat-ce-motivic-origin}
The Hodge structure on $H^\bullet_{\mathrm{CE}}(\mathfrak{g}_{\Delta_5})$
is entirely determined by the K3 period datum
\[
 \mathrm{Per} \colon H^2(\mathrm{K3}, \mathbb{Z}) \longrightarrow \mathbb{C}, \qquad
 \mathrm{Per}(\gamma) \;=\; \int_\gamma \omega_{\mathrm{K3}},
\]
composed with the Mukai extension $H^*(\mathrm{K3}, \mathbb{Z}) =
\mathrm{II}_{4,20}$ carrying the weight-0 Mukai Hodge structure
(bigraded $H^{(p,q)}_{\mathrm{Muk}}$ with $p+q = 0$; Mukai 1984
\emph{Nagoya Math.\ J.} 81). The K3 period $\omega_{\mathrm{K3}} \in
H^{2,0}(\mathrm{K3})$ is a well-defined element of $\mathbb{C}$ modulo
$\mathbb{Z}$-periods; its complex conjugate $\overline{\omega}_{K3}
\in H^{0,2}(\mathrm{K3})$ provides the anti-holomorphic input. The
resulting Hodge structure factors through the classifying space of
the Mukai period domain $\mathcal{D}_{\mathrm{Muk}} =
\mathrm{O}^+(\mathrm{II}_{4,20})\backslash \mathrm{O}(4,20)/\mathrm{SO}(2)\times
\mathrm{O}(2,20)$, whose Mumford--Tate group is the full orthogonal
group of the Mukai lattice.
\end{corollary}

\begin{remark}[Twisted period integrals and \texorpdfstring{$\mathbb{C}/\mathbb{Z}(k)$}{C/Z(k)}]
\label{rem:cycat-ce-twisted-periods}
\index{twisted period integral!$\mathfrak{g}_{\Delta_5}$ imaginary roots}
\index{motivic Tate twist!K3 period lift}
For each imaginary simple root $\alpha \in \Delta_+^{\mathrm{im}}$
with multiplicity $c_{\mathrm{K3}}(4nm-\ell^2)$, define the
\emph{twisted K3 period integral}
\[
 \pi_\alpha^{(k)} \;=\;
 \int_{\gamma_\alpha} \phi_{0,1}^{\mathrm{K3}} \cdot \omega_{\mathrm{K3}}^{k}
 \;\in\; \mathbb{C}/\mathbb{Z}(k),
\]
where $\gamma_\alpha \in H_2(\mathrm{K3}, \mathbb{Z})$ is the
$\alpha$-specific 2-cycle determined by the embedding $\alpha
\hookrightarrow H^2(\mathrm{K3}, \mathbb{Z})$ and $\mathbb{Z}(k) =
(2\pi i)^k \mathbb{Z}$ is the Tate twist. These twisted periods
are the entries of the motivic period matrix of the
$\alpha$-graded piece $\Lambda^k \mathfrak{g}_\alpha^*$ of CE
cohomology: $\mathrm{gr}^W_{2k} \mathrm{gr}^F_p H^k_{\mathrm{CE},\alpha}$
is generated over $\mathbb{Q}$ by $\pi_\alpha^{(p)} \cdot
(\overline{\pi_\alpha})^{k-p}$, modulo the symmetry
$\pi_\alpha^{(k)}|_{\omega \to \bar\omega}$. The $\mathbb{Z}(k)$
Tate twist reflects the $k$-fold wedge: the de Rham / Betti
comparison matrix on $H^k_{\mathrm{CE},\alpha}$ is a period matrix
with entries in $\mathbb{C}$ whose determinant is the
$\alpha$-indexed power of $(2\pi i)^{k}$, up to algebraic.
Primary: Deligne 1971 \emph{Inst.\ Hautes \'Etudes Sci.}~40 \S 2
(Hodge II mixed Hodge structures); Deligne 1974
\emph{Inst.\ Hautes \'Etudes Sci.}~44 \S 3 (Hodge III weight filtrations);
Bruinier LNM~1780 \S 2.2 (Borcherds singular-theta Hodge
decomposition).
\end{remark}

\begin{theorem}[Mixed Tate structure on CE cohomology; weight-\texorpdfstring{$n$}{n} from \texorpdfstring{$\zeta(n)$}{zeta(n)}]
\label{thm:cycat-ce-mixed-tate}
The weight-$n$ graded piece $\mathrm{gr}^W_n H^\bullet_{\mathrm{CE}}(\mathfrak{g}_{\Delta_5})$,
restricted to the real-root sector (pure Tate by
Theorem~\textup{\ref{thm:cycat-ce-hodge-real-imag}}\textup{(i)}) plus the
motivic piece of the imaginary-root sector that factors through
the $\mathbb{Z}(n)$-twist, is a \emph{mixed Tate motive} in the sense
of Deligne--Goncharov (\emph{Ann.\ Sci.\ ENS}~38, 2005) over $\mathbb{Z}$:
\begin{enumerate}[label=(\alph*)]
 \item The motivic Galois group of the sub-category generated by
 $\{\mathrm{gr}^W_n H^k_{\mathrm{CE}}(\mathfrak{g}_{\Delta_5})\}_n$ is
 $\mathbb{G}_m \ltimes U^{\mathrm{MT}}$ where $U^{\mathrm{MT}}$ is
 the pro-unipotent motivic Galois group, with $\mathrm{Lie}(U^{\mathrm{MT}})
 = \bigoplus_{n \geq 3, n\text{ odd}}\mathbb{Q}\cdot f_n$ (a free graded
 Lie algebra on $\{f_3, f_5, f_7, \ldots\}$ by Deligne--Goncharov).
 \item Period realisation of each $f_n$ generator is $\zeta(n)/
 (2\pi i)^n$ for odd $n \geq 3$. The \emph{weight-$n$ pieces of CE
 cohomology come from $\zeta(n)$-period pairings} between
 $\mathfrak{g}_\alpha^*$-classes in $H^1$ and co-classes in $H^{n-1}$.
 \item At weight 12, the first depth-4 MZV $\zeta(3,3,3,3)$ appears
 irreducibly, matching the Etingof weight-12 associator coefficient
 $\phi^{(12)}$ and entering the Borcherds product $\Delta_5^6
 \cdot (\text{weight-12 correction})$ at the same order.
\end{enumerate}
The mixed Tate structure is \emph{not} present in the Monster BKM
$\mathfrak{m}_{\mathrm{Monster}}$ (which has Cartan rank 2 and
pure Tate denominator $j(\sigma) - j(\tau)$); it arises in
$\mathfrak{g}_{\Delta_5}$ \emph{precisely} because the K3
holomorphic period $\omega_{\mathrm{K3}}$ enters the imaginary-root
multiplicities through the Hodge decomposition of $\phi_{0,1}^{\mathrm{K3}}$,
introducing period-theoretic transcendentals in the CE differential.
\end{theorem}

\begin{remark}[Motivic fundamental group of K3 and \texorpdfstring{$\mathrm{GRT}_1$}{GRT-1}-torsor]
\label{rem:cycat-ce-motivic-fundamental-group}
\index{motivic fundamental group!K3}
\index{Grothendieck--Teichm\"uller!super-GRT action on K3 periods}
The motivic fundamental group $\pi_1^{\mathrm{mot}}(\mathrm{K3})$ of
a complex algebraic K3 surface, after basepoint fixing and
unipotent completion, is \emph{non-trivial}: it contains a
pro-unipotent part coming from iterated period integrals of
$\phi_{0,1}^{\mathrm{K3}}$ against itself on $\mathrm{K3}$
(Brown--Levin 2011 iterated-integral formalism; Brown 2012
\emph{Ann.\ Math.}~175 depth theorem for periods of
$\pi_1^{\mathrm{mot}}(\mathbb{P}^1\setminus\{0,1,\infty\})$). The
resulting pro-unipotent motivic group receives a natural action of
the \emph{Grothendieck--Teichm\"uller group} $\mathrm{GRT}_1$
(Drinfeld 1990 \emph{Leningrad Math.\ J.}~2): the tangential-base-point
action of $\mathrm{Gal}(\overline{\mathbb{Q}}/\mathbb{Q})$ on K3 periods
factors through $\mathrm{GRT}_1$ by the Deligne--Ihara conjecture
(proved in the mixed Tate case by Brown 2012).

The space of Etingof--Kazhdan quantisations $\mathbf{H}_{\Delta_5}$
is a torsor over the super-Grothendieck--Teichm\"uller group
$\mathrm{GRT}_1^{\mathrm{super}}$, reflecting the $A_\infty$-quasi-Hopf
freedom in the associator $\widetilde\Phi^{\mathrm{Sieg-Bor}}$. The
\emph{same action}, via period-theoretic transfer, is realised on
the K3 motivic fundamental group: fixing a Drinfeld associator
$\Phi_{\mathrm{KZ}}$ on $\pi_1^{\mathrm{mot}}(\mathbb{P}^1\setminus
\{0,1,\infty\})$ determines, by Willwacher's theorem on the
bijection $\mathrm{GRT}_1 \simeq H^0(\mathrm{GC}_2)$ (graph
cohomology; Willwacher 2015 \emph{Invent.}~200), a class in $H^0$
of the Kontsevich graph complex, which pulls back to
$\pi_1^{\mathrm{mot}}(\mathrm{K3})$ via the K3 Kuga--Satake
embedding into a Shimura variety of orthogonal type. The
\emph{$\mathrm{GRT}_1$-torsor on $\mathbf{H}_{\Delta_5}$-quantisations
equals the $\mathrm{GRT}_1$ motivic Galois action on K3 periods},
and both act on $H^\bullet_{\mathrm{CE}}(\mathfrak{g}_{\Delta_5})$
compatibly with the Hodge filtration of
Theorem~\textup{\ref{thm:cycat-ce-hodge-real-imag}}.

Three comparison theorems enter:
\begin{itemize}
 \item \emph{Brown 2012}: the depth filtration on multiple zeta values
 matches the Mukai-lattice rank grading on imaginary-root
 multiplicities (both bounded by 4 at weight 12); the first
 depth-4 class $\zeta(3,3,3,3)$ enters both
 $\mathrm{Lie}(U^{\mathrm{MT}})$ and the weight-12 CE cohomology.
 \item \emph{Willwacher 2015}: $\mathrm{GRT}_1 \simeq H^0(\mathrm{GC}_2)$
 implies the K3-Kuga--Satake graph-cohomology transfer, tensorially
 compatible with Etingof--Kazhdan's super-GRT parametrisation.
 \item \emph{Deligne 1989 }\emph{Galois groups over $\mathbb{Q}$}: the mixed
 Tate motivic Galois group over $\mathbb{Z}$ surjects onto the K3
 motivic Galois sub-quotient, with kernel controlled by the
 non-Tate pieces (the imaginary-root mixing of
 Theorem~\textup{\ref{thm:cycat-ce-hodge-real-imag}}\textup{(ii)}).
\end{itemize}
Primary: Brown 2012 \emph{Ann.\ Math.}~175; Willwacher 2015
\emph{Invent.}~200; Deligne--Goncharov 2005 \emph{Ann.\ Sci.\ ENS}~38;
Kontsevich 1999 \emph{Lett.\ Math.\ Phys.}~48; Mukai 1984
\emph{Nagoya Math.\ J.}~81; Etingof--Kazhdan 2000 Part V
($\phi^{(n)}$-polynomial MZV content in the super setting).
\end{remark}
